% ============================================================================
% Recovering a Taxonomy of Collective Behaviors from Motion Alone
%
% Set with the Springer Nature template (sn-jnl.cls, unmodified) in the APA reference style that the journal Swarm
% Intelligence uses (option sn-apa, sn-apacite.bst). This file is the arXiv version: the paper with its appendices.
% journal.tex compiles the same text without the appendices, which supplementary.tex turns into the Supplementary
% Information of the journal version. The shared preamble is preamble.tex; the related work and the appendices are
% input from sections/.
% ============================================================================
\documentclass[pdflatex,sn-apa]{sn-jnl}
% ============================================================================
% Preamble shared by the three documents built from one text:
%   paper.tex         arXiv version: the paper with every appendix in the same PDF;
%   journal.tex       Swarm Intelligence version: the paper alone, the appendices going to
%   supplementary.tex its Supplementary Information, numbered S1, S2, ...
% The two journal documents read each other's labels with xr-hyper (loaded before the class).
% Every number comes from numbers.tex, generated by analysis/make_numbers.py from the result files, which also
% checks the statements of the text; tables come from tables/. Missing values print in red. The figures are drawn
% at their printed width in the journal (122 mm), with Arial or Helvetica lettering of at least 8.5 pt.
% ============================================================================
\usepackage[T1]{fontenc}
\usepackage{lmodern}
\usepackage{amsmath,amssymb,mathtools}
\usepackage{graphicx}
\usepackage{booktabs}
\usepackage{multirow}
\usepackage{array}
\usepackage{tabularx}
\usepackage{xcolor}
\usepackage{tikz}
\usepackage{placeins}
\usepackage{pifont}
\newcommand{\cmark}{\ding{51}}
% floats: up to two at the top of a page, and pages of floats when needed, so that none drifts far from its text
\renewcommand{\topfraction}{0.9}\renewcommand{\bottomfraction}{0.6}\renewcommand{\textfraction}{0.08}
\renewcommand{\floatpagefraction}{0.75}\setcounter{topnumber}{3}\setcounter{bottomnumber}{1}\setcounter{totalnumber}{4}
\usepackage[capitalise,noabbrev]{cleveref}
% references in the style of the journal: Fig. 1, Table 1, Section 2, Eq. (1), Appendix A
\crefname{figure}{Fig.}{Figs.}\Crefname{figure}{Fig.}{Figs.}
\crefname{equation}{Eq.}{Eqs.}\Crefname{equation}{Eq.}{Eqs.}
\crefname{proposition}{Proposition}{Propositions}\Crefname{proposition}{Proposition}{Propositions}
\crefname{appendix}{Appendix}{Appendices}\Crefname{appendix}{Appendix}{Appendices}
\newcommand{\doilink}[1]{\href{https://doi.org/#1}{https://doi.org/#1}}
% APA 7 reference list, as the journal prints it: an ampersand before the last of three or more authors
% (the Springer APA style writes it only for two authors)
\usepackage{etoolbox}
% tables: 9 pt, a little air between the rows, the row of the method shaded
\usepackage{makecell}
\usepackage{colortbl}
\newcommand{\methodrow}{\rowcolor{black!7}}
\AtBeginEnvironment{table}{\small\renewcommand{\arraystretch}{1.15}}
\AtBeginEnvironment{APACrefauthors}{\renewcommand{\BCBL}{,\ \BBA}}
% and no "Retrieved from" before a URL, which APA 7 dropped
\AtBeginDocument{\renewcommand{\BRetrievedFrom}{}}
% figures and tables of the appendices are numbered within each appendix: Fig. C1, Table B2
\AtBeginEnvironment{appendices}{\counterwithin*{figure}{section}\counterwithin*{table}{section}%
  \crefalias{section}{appendix}}
% the ORCID logo of the class is an EPS file that the template does not ship: the same mark is drawn here
\definecolor{orcidgreen}{HTML}{A6CE39}
\renewcommand{\orcidlogo}{\tikz[baseline=-0.55ex]{\fill[orcidgreen] (0,0) circle (0.62em);
  \node[white, font=\sffamily\bfseries\fontsize{6.5}{6.5}\selectfont, inner sep=0pt] at (0,0) {iD};}}
% clean page breaks: never one or two lines of a paragraph alone at the bottom or the top of a page
\clubpenalty=10000 \widowpenalty=10000 \displaywidowpenalty=10000 \brokenpenalty=10000
\setlength{\emergencystretch}{3em}
% paragraphs: an indent and a small skip between them
\setlength{\parindent}{1.2em}
\setlength{\parskip}{0.5ex}
% regular space around figures and tables: fixed gaps, and pages that are not stretched to the bottom, so that no
% gap above or below a float grows to fill a page
\raggedbottom
\setlength{\textfloatsep}{14pt plus 1pt minus 1pt}\setlength{\floatsep}{12pt plus 1pt minus 1pt}
\setlength{\intextsep}{12pt plus 1pt minus 1pt}
\setlength{\abovecaptionskip}{6pt}\setlength{\belowcaptionskip}{6pt}

\theoremstyle{thmstyleone}
\newtheorem{proposition}{Proposition}

% ---------------------------------------------------------------- notation
\newcommand{\R}{\mathbb{R}}
\newcommand{\knn}{k_{\mathrm{nn}}}
\newcommand{\vbar}{\bar v}
\newcommand{\Sstar}{S^{\star}}
\newcommand{\Cstar}{C^{\star}}
\newcommand{\SW}{\mathrm{SW}}
% vectors carry an arrow, scalars none, and the second moment, a tensor, is in bold
\newcommand{\vx}{\vec{x}}
\newcommand{\vv}{\vec{v}}
\newcommand{\vr}{\vec{r}}
\newcommand{\vd}{\vec{d}}
\newcommand{\vV}{\vec{V}}
\newcommand{\vM}{\vec{M}}
\newcommand{\vW}{\vec{W}}
\newcommand{\vVb}{\overline{\vV}}
\newcommand{\tM}{\mathbf{M}}
% numbers generated from the results (manuscript/numbers.tex)
\newcommand{\setnum}[2]{\expandafter\def\csname num@#1\endcsname{#2}}
\newcommand{\nm}[1]{\ifcsname num@#1\endcsname\csname num@#1\endcsname\else\textcolor{red}{\textbf{[#1]}}\fi}
\InputIfFileExists{numbers.tex}{}{}
\newcommand{\tabin}[1]{\InputIfFileExists{tables/#1.tex}{}{\textcolor{red}{\textbf{[table \detokenize{#1} pending]}}}}
\newcommand{\figin}[2]{\IfFileExists{#2}{\includegraphics[width=#1]{#2}}%
{\fbox{\parbox[c][2cm][c]{#1}{\centering\small\textcolor{red}{figure pending}\\\detokenize{#2}}}}}
% where the material that is not in the body of the paper lies: the appendices (arXiv) or the Supplementary
% Information (journal)
\ifdefined\journalversion
  \newcommand{\suppl}{the Supplementary Information}
\else
  \newcommand{\suppl}{the appendices}
\fi
% the full list of cited works, so that APA 7 prints "et al." from the first citation of three or more authors
\newcommand{\allshortcites}{\shortcites{vicsek1995novel,vicsek2012collective,couzin2002collective,dorsogna2006self,%
gomes2014systematic,berndt1994using,topaz2015topological,corcoran2017modelling,bhaskar2019analyzing,xian2022capturing,%
zeng2024neural,suetani2025unsupervised,legarda2026universal,edelsbrunner2010computational,carlsson2009topology,%
edelsbrunner2002topological,zomorodian2005computing,cohen2009extending,cohen2007stability,maria2014gudhi,%
carriere2017sliced,von2007tutorial,zelnik2004self,shi2000normalized,bhatia2013helmholtz,pinciroli2012argos,%
mondada2009puck,francesca2014automode,francesca2015chocolate,brambilla2013swarm,uci_swarm,wang2014similarity,%
obayashi2009rank,frey2007clustering,benhur2002stability,campello2013density,edelsbrunner2004jacobi,%
hasselmann2021empirical,gomes2013generic,fialho2026differ,mattson2023exploring,brown2018discovery,%
tunstrom2013collective,sathiyakumar2024detecting,ndinga2025thesis,monti2003consensus,rosenberg2007vmeasure,%
hubert1985comparing,carlsson2009theory,botnan2023introduction,leygonie2022framework,carriere2021optimizing,%
bhatia2014natural,tong2003discrete,petronetto2010meshless,weinhart2012discrete,schranz2020swarm,rossi2018review,%
mattson2023leveraging,mattson2025discovery,titus2021unsupervised,alhubaishi2026interpretable,abpeikar2023motion,%
li2016turing,berman2014mapping,berman2016predictability,klibaite2017unsupervised,mototake2024topological,%
kern2026study,wang2025deep,polthier2003identifying,macedo2010learning,witkin1983scale,lindeberg1994scale}}

\ifdefined\journalversion\externaldocument[][nocite]{supplementary}[]\fi

\begin{document}
\allshortcites

\title[Recovering a Taxonomy of Collective Behaviors from Motion Alone]{Recovering a Taxonomy of Collective Behaviors
from Motion Alone}

\author*[1]{\fnm{Olivier Vertu} \sur{Ndinga Oba}\,\orcid{https://orcid.org/0009-0001-3422-2846}}\email{olivier.ndinga.oba@ulb.be}

\affil*[1]{\orgname{Universit\'e libre de Bruxelles}, \orgaddress{\city{Brussels}, \country{Belgium}}}

\abstract{Tracking a robot swarm or animal group gives the trajectories of its agents, but not which collective behaviors
those recordings perform. From tracked positions alone, we recover an unsupervised taxonomy of the behaviors: their number,
their members, one exemplar each, their resolution hierarchy, and what in the motion separates them. A Gaussian footprint
turns each recording into five rotation-invariant fields at four data-driven scales: local density, kinetic intensity,
the two potentials of a closed-form Helmholtz--Hodge decomposition of the flux, and their Jacobi coupling. Extended
persistence summarizes each field without a grid. Rank-normalized sliced Wasserstein distances are fused across scales
without weights fitted to labels, and affinity propagation returns the behaviors and their exemplars. Across
\nm{n-runs} simulated runs of 20 e-puck robots in five missions and up to ten controllers per mission, one controller
per mission (\nm{top1-n} runs) gives a modal consensus (\nm{top1-plateau} of 99 settings) that places every run with its
mission; \nm{top1-errfree} settings misassign no run, and \nm{top1-errfree-plateau} recover the five missions exactly.
With up to ten controllers per mission (up to \nm{n-runs} runs), fine resolutions split missions into pure
sub-strategies (purity at least \nm{pooled-min-purity}), intermediate ones recover the five missions, and coarse ones
merge traffic (Directional Gate with Foraging) and gathering (Shelter with XOR Aggregation), while Homing stays apart.
Because the distance splits over critical structures and is differentiable, differences between behaviors trace back
to rotation centers or sinks and to the robots that shaped them.}

\keywords{Swarm robotics, Topological data analysis, Similarity network fusion, Affinity propagation, Unsupervised
behavior discovery}

\maketitle

% ============================================================================
\section{Introduction}
\label{sec:intro}

Collective behaviors are among the most striking forms of coordination: starlings fold into murmurations, fish schools
switch between milling and polarized swimming \citep{tunstrom2013collective}, and swarms of simple robots aggregate or
forage without central control \citep{brambilla2013swarm}. When a robot swarm is designed
automatically, different controllers can even solve the same mission with different motion strategies. Recordings of
collectives rarely state which behaviors they contain, so comparing, designing or imitating collectives requires a way
to identify the behaviors from the motion alone.

We call a collective behavior the organization of the motion of the agents that recurs across recordings: where they
gather and disperse, how they circulate and rotate, from each agent to the whole group. Two recordings share a
behavior when their motion is organized alike, whatever controller or task produced it, and \cref{sec:method} turns
this likeness into a distance. Our goal is a taxonomy of these behaviors. In unsupervised learning, a taxonomy can be discovered from unlabeled data
as a hierarchy of categories, each represented by prototypes, without assuming their number
\citep{wang2025deep}. Here, with no label and no number of behaviors given, a taxonomy is the data of (i)~the
number of behaviors, (ii)~the partition of the recordings, (iii)~one exemplar recording per behavior, (iv)~the
hierarchy of the behaviors across resolutions, from groups within a behavior to families of behaviors close in motion,
and (v)~the attribution of their differences to structures of the motion and to the agents that shaped them. Existing
approaches return some of these components, but not all of them (\cref{sec:related}).

Someone who watches a school of fish sees each fish and the whole school at once. Our central hypothesis is that a
collective must be read in the same way, at several scales at once, which fields make possible: a Gaussian footprint
spreads each agent over a width that acts as the observation scale. Our previous work used a topological similarity of velocity and
density fields to imitate swarm behaviors \citep{ndinga2025thesis}; here, fields and topology serve a taxonomy.

This paper makes four contributions, summarized in \cref{fig:overview}. (i)~Each recording is represented by five
rotation-invariant fields at four scales derived from the data: the density, the kinetic intensity, the two
Helmholtz--Hodge potentials of the flux, in closed form on the whole plane, and their Jacobi coupling. (ii)~Each field
is summarized by its extended persistence, computed without a grid. Sliced Wasserstein distances, rank-normalized and
fused across scales by similarity network fusion (SNF) without weights tuned on labels, lead to affinity propagation
(AP), which returns the number of behaviors, the partition and the exemplars. (iii)~Sweeping the resolution of the
fusion gives the hierarchy. (iv)~The distance splits exactly over the critical structures of the fields and is
differentiated back to the agents. On \nm{n-runs} simulated runs of $20$ e-puck robots in five missions, the method
recovers the five missions without labels, reveals pure sub-strategies among pooled controllers, merges the missions
that move alike at coarse resolution, and traces their differences back to the rotation centers of their traffic.

% ============================================================================
\section{Related work}
\label{sec:related}

The taxonomies of swarm robotics are built a priori. \citet{brambilla2013swarm} classified the field by its design and
analysis methods and by collective behavior, into spatially organizing behaviors, navigation, collective
decision-making and others. \citet{schranz2020swarm} applied such classes of basic behaviors to current applications,
and \citet{rossi2018review} classified the algorithms by their mathematical structure. These taxonomies organize what a
swarm can be designed to do, not what a set of recordings contains.

Behaviors are usually compared through descriptors chosen in advance: order parameters such as polarization and
milling \citep{vicsek1995novel,couzin2002collective,dorsogna2006self}, which reveal the swarm, milling and polarized
states of fish schools \citep{tunstrom2013collective}, or the behavior characterizations of swarm robotics
\citep{gomes2013generic,gomes2014systematic}. Which groups a comparison finds depends on the pairing of these features
with a similarity measure \citep{fialho2026differ}. Classifiers learn instead human labels of flocking, alignment and
grouping \citep{abpeikar2023motion}, or the labeled phases of a real robot swarm, with which unsupervised clusterings
agree only moderately \citep{alhubaishi2026interpretable}.

Other works let behaviors emerge. Novelty search explores the controllers of limited robots, and k-medoids clusters
the archive into a taxonomy whose medoids are representative behaviors, their number chosen by inspection
\citep{brown2018discovery}. The behavior space was later learned with human feedback, and the search
extended to heterogeneous swarms and real robots \citep{mattson2023leveraging,mattson2023exploring,mattson2025discovery}.
\citet{titus2021unsupervised} compare by diffusion maps the networks that positions and velocities define over a group,
and cluster the result into three groups fixed in advance. Computational ethology goes furthest toward a taxonomy
discovered from motion, but for single animals: the postural dynamics of flies divide, by a watershed of their density,
into more than one hundred stereotyped behaviors organized hierarchically \citep{berman2014mapping,berman2016predictability},
and the approach extends to pairs of animals \citep{klibaite2017unsupervised}.

Topological data analysis summarizes shapes instead of chosen descriptors. \citet{topaz2015topological} summarized in
CROCKER plots the persistent homology of the point clouds that agents form in the space of positions and velocities. On
positions alone, such summaries classify the regimes of the D'Orsogna model better than order parameters
\citep{bhaskar2019analyzing} and separate the regimes of a crowd \citep{kern2026study}. \citet{corcoran2017modelling}
followed a school of fish by zig-zag persistence of a density estimated on a grid. Other works follow these summaries in
time \citep{xian2022capturing,sathiyakumar2024detecting}, embed persistence images \citep{suetani2025unsupervised},
learn latent dynamics of diagrams \citep{zeng2024neural}, or measure the fluctuations of a large flock across
hierarchical levels \citep{mototake2024topological}. In all of them, the topology is that of a point cloud or of a
density, and velocities, when used, are coordinates of points rather than a field. \citet{legarda2026universal}
described instead the topology of a velocity field, through the extended persistence of its Helmholtz--Hodge
potentials, and grouped the recordings into a fixed number of clusters or by HDBSCAN. \Cref{tab:related} summarizes what
each of these approaches provides.

\begin{table}[!htbp]
\centering
\caption{What each approach provides (\cmark\ yes, -- no). \emph{medoids}: k-medoids for a number of clusters chosen
by the user; \emph{partly}: with one of the clustering methods used; \emph{features}: importances of hand-crafted
features; \emph{time}: time scales; \emph{levels}: the levels of one flock.}
\label{tab:related}
\newcommand{\appr}[2]{#1\newline{\footnotesize\color{black!72}\citealp{#2}}}
\newcommand{\tw}[1]{{\footnotesize\emph{#1}}}
\newcommand{\rot}[1]{\multicolumn{1}{c}{\rotatebox[origin=l]{90}{#1}}}
\setlength{\tabcolsep}{1pt}
\renewcommand{\arraystretch}{1.0}
\newcommand{\rowsep}{\arrayrulecolor{black!22}\specialrule{0.4pt}{1.5pt}{1.5pt}\arrayrulecolor{black}}
\newcommand{\groupsep}{\arrayrulecolor{black!55}\specialrule{0.6pt}{2pt}{1.5pt}\arrayrulecolor{black}}
\newcolumntype{C}[1]{>{\centering\arraybackslash}p{#1}}
% the approach column takes what the eight narrow columns (6.81 cm) and their padding (16 x 1 pt) leave of the text
% width, so that the table is exactly as wide as the text and the shaded row is continuous
\begin{tabular}{@{}>{\raggedright\arraybackslash}p{\dimexpr\linewidth-6.81cm-16pt\relax}C{0.72cm}C{0.82cm}C{0.62cm}C{0.8cm}C{0.8cm}C{1.0cm}C{1.05cm}C{1.0cm}@{}}
\toprule
& \multicolumn{3}{c}{Representation} & \multicolumn{5}{c}{Taxonomy returned}\\
\cmidrule(lr){2-4}\cmidrule(l){5-9}
Approach & \rot{Collective} & \rot{Velocities} & \rot{Scales} & \rot{No labels} & \rot{Number} & \rot{Exemplar} &
\rot{Hierarchy} & \rot{Attribution}\\
\midrule
\multicolumn{9}{@{}l}{\emph{Taxonomies set a priori}}\\
\appr{Classes of behaviors or algorithms}{brambilla2013swarm,schranz2020swarm,rossi2018review}
  & \cmark & -- & -- & -- & -- & -- & \tw{a priori} & --\\
\groupsep
\multicolumn{9}{@{}l}{\emph{Descriptors chosen in advance}}\\
\appr{Order parameters}{vicsek1995novel,couzin2002collective,tunstrom2013collective}
  & \cmark & \cmark & -- & -- & -- & -- & -- & --\\
\rowsep
\appr{Behavior characterizations}{gomes2013generic,fialho2026differ}
  & \cmark & \tw{varies} & -- & -- & -- & -- & -- & --\\
\rowsep
\appr{Supervised classifiers and clustering}{abpeikar2023motion,alhubaishi2026interpretable}
  & \cmark & \cmark & -- & \tw{partly} & \tw{partly} & -- & -- & \tw{features}\\
\groupsep
\multicolumn{9}{@{}l}{\emph{Behaviors discovered without labels}}\\
\appr{Novelty search, then k-medoids}{brown2018discovery,mattson2023leveraging,mattson2023exploring,mattson2025discovery}
  & \cmark & \tw{varies} & -- & \cmark & -- & \tw{medoids} & -- & --\\
\rowsep
\appr{Manifold learning}{titus2021unsupervised}
  & \cmark & \cmark & -- & \cmark & -- & -- & -- & --\\
\rowsep
\appr{Behavioral maps of animals}{berman2014mapping,berman2016predictability,klibaite2017unsupervised}
  & \tw{pairs} & \cmark & \tw{time} & \cmark & \cmark & -- & \cmark & --\\
\groupsep
\multicolumn{9}{@{}l}{\emph{Topology of positions}}\\
\appr{CROCKER plots}{topaz2015topological,bhaskar2019analyzing,kern2026study}
  & \cmark & \tw{partly} & \cmark & \cmark & -- & \tw{medoids} & -- & --\\
\rowsep
\appr{Zig-zag persistence}{corcoran2017modelling}
  & \cmark & -- & -- & \cmark & -- & \tw{medoids} & -- & --\\
\rowsep
\appr{Persistence over time}{sathiyakumar2024detecting}
  & \cmark & -- & \cmark & -- & -- & -- & -- & --\\
\rowsep
\appr{Levels of a large flock}{mototake2024topological}
  & \cmark & -- & \cmark & -- & -- & -- & \tw{levels} & --\\
\groupsep
\multicolumn{9}{@{}l}{\emph{Topology of motion}}\\
\appr{Persistence of a velocity field}{legarda2026universal}
  & \cmark & \cmark & -- & \cmark & \tw{partly} & -- & -- & --\\
\groupsep
\methodrow
\textbf{This work} & \cmark & \cmark & \cmark & \cmark & \cmark & \cmark & \cmark & \cmark\\
\bottomrule
\end{tabular}
\end{table}

Each of these lines of work addresses its own problem, from organizing the design of swarms to characterizing physical
regimes, mapping the postures of single animals or comparing velocity fields, and each provides part of what a
taxonomy of collective behaviors requires. What remains open is to join a topological description of collective
motion with the unsupervised discovery of a taxonomy and with an explanation of its differences down to the agents,
which is what \cref{sec:method} presents.


\FloatBarrier
% ============================================================================
\section{Methodology}
\label{sec:method}

We are given many recordings of identical agents that form a swarm. A recording is the set of their positions over time,
$\{\vx_i(t)\}_{i=1}^N\subset\R^2$, sampled every $\Delta t$ at $T$ instants. These positions are the only input of the
method, and labels are used only to score its result. The method rests on one idea: a collective must be read at
several scales at once, because the agents and the group they form are both part of the behavior. We therefore spread
the agents into continuous fields with a Gaussian footprint and read these fields at several widths
(\cref{sec:fields}), before turning them into a taxonomy (\cref{sec:from_fields}). \Cref{fig:overview} follows the whole
pipeline on the robot data, and \cref{app:method} details and illustrates each step.

\begin{figure}[!htbp]
\centering
\figin{\linewidth}{figures/fig_method.pdf}
\caption{The method on the robot data. Top: one recording, the Shelter exemplar, from its tracks to its diagrams;
in (c), seeds at the robots (gray), near critical points (orange) and on $J_\sigma=0$ (blue). Bottom, from right to
left: the $250$ runs of Top-1, from their distances to their $C=5$ behaviors. Dashed: the attribution.}
\label{fig:overview}
\end{figure}

\subsection{From motion to invariant fields}
\label{sec:fields}\label{sec:representation}\label{sec:ladder}

Newton measured the motion of a body by its quantity of motion, which arises from its velocity and its quantity of
matter together, the motion of a whole being the sum of the motions of its parts
\citep[Definition~II]{newton1687principia}. For identical agents, the motion of each is its velocity, the forward
difference $\vv_i(t)=\big(\vx_i(t+\Delta t)-\vx_i(t)\big)/\Delta t$. The motion of a collective lies in how its agents
are distributed over the positions and the velocities. Its description must also be unchanged by the rotations of the
plane, the group $\mathrm{SO}(2)$, and by its translations, since a group moving north performs the same behavior as a
group moving east.

We spread each agent over a Gaussian footprint of width $\sigma$ and take the first three moments of the velocity
distribution weighted by it:
\begin{equation}
O_i(\vr,t)=\exp\!\Big(-\frac{\|\vr-\vx_i(t)\|^2}{2\sigma^2}\Big),
\label{eq:footprint}
\end{equation}
\begin{equation}
\begin{gathered}
M_0(\vr,t)=\sum_{i=1}^N O_i(\vr,t),\qquad
\vM_1(\vr,t)=\sum_{i=1}^N \vv_i(t)\,O_i(\vr,t),\\
\tM_2(\vr,t)=\sum_{i=1}^N \vv_i(t)\otimes\vv_i(t)\,O_i(\vr,t).
\end{gathered}
\label{eq:moments}
\end{equation}
For agents of unit mass, $M_0$ is the density and $\vM_1$ the momentum density, the fields by which particle systems
are coarse-grained into continua \citep{weinhart2012discrete}, and $\operatorname{tr}\tM_2$ is twice the kinetic
energy density. The footprint favors no direction and is negligible beyond $5\sigma$ (\cref{prop:so2,prop:tail}), it
is smooth in the positions, and it gives the potentials below in closed form (\cref{prop:hhd}). Its width acts as the observation scale: at a small width every agent is
resolved, and at a large width the agents merge into the groups they form.

Averaged over the $T$ instants, which keeps the organization of the traffic and leaves out its chronology, the moments
give the density and the kinetic intensity,
\begin{equation}
\rho_\sigma(\vr)=\frac{1}{NT}\sum_t M_0(\vr,t),\qquad
K_\sigma(\vr)=\frac{1}{NT\,\overline{v^2}}\sum_t\operatorname{tr}\tM_2(\vr,t),
\label{eq:density}
\end{equation}
where the mean squared speed $\overline{v^2}$ of the recording makes $K_\sigma$ independent of the rate at which the
recording is replayed. Both fields are local, evaluated within $5\sigma$ of the agents. Where two streams cross in
opposite directions, they cancel in the net flux below but add up in $K_\sigma$ (\cref{fig:moments}). The flux $\vV_\sigma=\vM_1$,
divided by $\sigma\vbar$ with $\vbar$ the mean speed, is split by the Helmholtz--Hodge decomposition
\citep{bhatia2013helmholtz},
\begin{equation}
\vVb_\sigma(\vr)=\frac{1}{NT\sigma\vbar}\sum_t\vV_\sigma(\vr,t)=\nabla\varphi_\sigma+\nabla^{\perp}\psi_\sigma,\qquad
\nabla^{\perp}=(-\partial_y,\partial_x).
\label{eq:hhd}
\end{equation}
The divergence potential $\varphi_\sigma$ captures the sinks and sources of the flux, its maxima and minima, and the
vorticity potential $\psi_\sigma$, a stream function, its centers of rotation, its extrema. Because the flux is a sum of
Gaussians on the whole plane, both potentials have a closed form through a single radial kernel,
\begin{equation}
\kappa_\sigma(\vd)=\frac{\sigma^2}{\|\vd\|^2}\Big(1-e^{-\|\vd\|^2/(2\sigma^2)}\Big),
\label{eq:kappa}
\end{equation}
which equals $1/2$ at an agent and decays as $1/\|\vd\|^2$, so that the potentials, unlike the density, also depend on
distant agents. With $\vd_i(t)=\vr-\vx_i(t)$, $u_i=\vv_i\cdot\vd_i$ and $w_i=v_{i,x}d_{i,y}-v_{i,y}d_{i,x}$,
\begin{equation}
\varphi_\sigma(\vr)=\frac{1}{NT\sigma\vbar}\sum_{t,i}u_i(t)\,\kappa_\sigma\big(\vd_i(t)\big),\qquad
\psi_\sigma(\vr)=-\frac{1}{NT\sigma\vbar}\sum_{t,i}w_i(t)\,\kappa_\sigma\big(\vd_i(t)\big).
\label{eq:potentials}
\end{equation}
These potentials rebuild the averaged flux exactly, without a grid, a boundary condition or a harmonic part
(\cref{prop:hhd}).

No single width can be the right one, so every field is computed on a single global ladder of $M$ widths read once from
the trajectories across every instant of all recordings, without labels. It runs from half the smallest distance
$d_{\min}$ between two agents, floored by the tracking precision $\epsilon_\sigma$, to a factor $k$ below the width
$D_s/2.35$ at which the footprints of the whole swarm, of largest diameter $D_s$, merge:
\begin{equation}
\sigma_{\min}=\max\!\Big(\frac{d_{\min}}{2},\epsilon_\sigma\Big),\quad
\sigma_{\max}=\frac{D_s}{2.35\,k},\quad
\sigma_m=\sigma_{\min}\Big(\frac{\sigma_{\max}}{\sigma_{\min}}\Big)^{\frac{m-1}{M-1}},\quad m=1,\dots,M.
\label{eq:ladder}
\end{equation}
Here $2.35\approx2\sqrt{2\ln2}$ is the full width at half maximum of a Gaussian of unit width: at $\sigma=D_s/2.35$, one
footprint spans the swarm at half its height. With $M=4$ widths and the cap $k=\nm{cap-k}$, the ladder runs from $\nm{sigma1-cm}$ to $\nm{sigma4-cm}$\,cm on the robot
data, and the taxonomy hardly depends on the cap (\cref{app:gallery}).

\subsection{From fields to a taxonomy}
\label{sec:from_fields}\label{sec:seeds}\label{sec:distances}\label{sec:fusion}\label{sec:ap}\label{sec:differentiability}

At every width, each field is summarized by what it organizes rather than where: its sinks, sources, rotation centers
and the saddles between them, their values and the way they nest. Extended persistence provides this summary
\citep{cohen2009extending}. It follows the sublevel sets $\{f\le c\}$ of the field as $c$ rises, completed by its
superlevel sets swept downward, so that every component and every cycle has a finite birth and death. It yields six
diagrams per field, which change little when the field changes little. On a planar domain, $\mathrm{Ord}_0$ records how
the wells of the field, its minima inside the domain, merge at saddles as the level rises, and $\mathrm{Rel}_1$ how its
peaks, its maxima, merge at saddles as the level falls. Each interior peak is enclosed by a sublevel loop killed at the peak in $\mathrm{Ord}_1$, and each
interior well forms a relative $2$-cycle spanning from its minimum to its surrounding rim in $\mathrm{Rel}_2$.
$\mathrm{Ext}_0$ spans the field from its lowest to its highest value, and $\mathrm{Ext}_1$ records the holes
of the domain (\cref{app:method}). For the potentials, the peaks are sinks or clockwise rotation centers, and the wells
sources or counterclockwise ones. Taken separately, the diagrams of the two
potentials do not say whether a sink and a rotation center lie at the same place. A fifth field records it, the
alignment of their gradients,
\begin{equation}
J_\sigma(\vr)=\sigma^2\det\big(\nabla\varphi_\sigma,\nabla\psi_\sigma\big)(\vr),
\label{eq:jacobi}
\end{equation}
which vanishes on the Jacobi set of the pair \citep{edelsbrunner2004jacobi} and is large where a sink and a rotation
center are apart (\cref{fig:potentials}). The five fields $\rho_\sigma$, $K_\sigma$, $\varphi_\sigma$, $\psi_\sigma$
and $J_\sigma$ are invariant under rotations and translations, and differentiable in the trajectories.

\begin{figure}[!htbp]
\centering
\figin{\linewidth}{figures/fig_potentials.pdf}
\caption{Two synthetic swarms with the same ingredients, $40$ agents gathering and $40$ milling ($\sigma=7$\,cm,
$30$\,s): when the sink and the rotation center coincide (top), $J_\sigma$ vanishes; when they are apart (bottom), it
is large between them. The diagrams of $\varphi_\sigma$ and $\psi_\sigma$ change by $\nm{cpl-phi-pct}$ and
$\nm{cpl-psi-pct}\,\%$ of their distance to the empty diagram.}
\label{fig:potentials}
\end{figure}

The diagrams are computed without a grid, whose spacing would set a chosen resolution. Each
recording gets its own seeds at each width: agent positions, points near the critical points of each field, and points
on the zero set of $J_\sigma$. The domain is the alpha complex of these seeds
\citep{edelsbrunner2010computational,maria2014gudhi}, with the lower-star filtration of each field
(\cref{app:method}).

Two recordings are compared through their diagrams by the sliced Wasserstein distance of \citet{carriere2017sliced}.
Each diagram is augmented with the projections of the points of the other onto the diagonal, and the two are projected
and sorted along $\nm{n-dirs}$ directions $\theta\in\Theta$ of a half-turn. With $a^\theta_{(k)}$ and $b^\theta_{(k)}$
the sorted projections, and $D^s_{f,m}(A)$ the six diagrams of recording $A$ on field $f$ at width $\sigma_m$,
\begin{equation}
\SW(D,D')=\frac{1}{|\Theta|}\sum_{\theta\in\Theta}\sum_{k}\big|a^{\theta}_{(k)}-b^{\theta}_{(k)}\big|,\qquad
d_{f,m}(A,B)=\sum_{s=1}^{6}\SW\big(D^{s}_{f,m}(A),D^{s}_{f,m}(B)\big).
\label{eq:sw}
\end{equation}
The distance is of order one, so that each structure counts in proportion to its persistence, and it is differentiable
almost everywhere. Over all pairs of recordings, it gives $5M=20$ matrices, one per field and width.

These matrices are combined without any weight tuned on labels. Each pair $(A,B)$ first receives its mutual rank in
$d_{f,m}$, the geometric mean of the rank of $B$ among the neighbors of $A$ and of the rank of $A$ among those of $B$
\citep{obayashi2009rank}, replaced by its percentile rank $\tilde d_{f,m}\in(0,1]$. Each moment of \eqref{eq:moments}
then receives the same weight, the flux sharing its weight among its three fields, which gives one view per width:
\begin{equation}
W_m=1-\Big(\tfrac13\tilde d_{\rho,m}+\tfrac13\tilde d_{K,m}+\tfrac19\big(\tilde d_{\varphi,m}+\tilde d_{\psi,m}
+\tilde d_{J,m}\big)\Big).
\label{eq:view}
\end{equation}
Similarity network fusion \citep{wang2014similarity} reads the four views together, so that a link supported by
several widths is kept. It keeps the $\knn$ largest affinities of each row of each view, normalized to sum to one, as
the local kernel $S_m$, and diffuses each row-normalized view $P_m$ through the mean of the others for $t$ iterations,
\begin{equation}
P_m\leftarrow S_m\Big(\frac{1}{M-1}\sum_{m'\neq m}P_{m'}\Big)S_m^{\top},
\label{eq:snf}
\end{equation}
each iteration being followed by a symmetrization and a row normalization. The consensus network $\Sstar$ is the
symmetrized mean of the diffused views.

Affinity propagation \citep{frey2007clustering}, applied to $\Sstar$ with the settings of its original implementation,
returns at once the number $C$ of behaviors, the partition and one exemplar per cluster, an actual recording
(\cref{app:method}). The result depends on the resolution $(\knn,t)$: small neighborhoods and short diffusion split the
recordings into pure sub-groups, and large ones merge them. Over a wide window whose neighborhoods grow proportionally
with the number of recordings, the partitions form nested \emph{stable states}, broad regions of settings that group the
recordings alike up to a few crossing runs, separated by \emph{transitions} where whole groups join. These words echo phase
diagrams, but $\knn$ and $t$ set the reading: the states belong to the resolution map, not to the swarm. Because coarse
states fill most of that wide window on pooled datasets, selecting the mode there would collapse behaviors into broad
families. We therefore evaluate the taxonomy over a fixed local grid of $99$ settings probing the fine scale. The
\emph{modal region} is the subset of this grid where $C$ takes its most frequent value $\Cstar$, and the taxonomy is its
\emph{modal consensus}: two recordings are linked by the fraction of modal settings in which they share a cluster, and
this matrix is cut into $\Cstar$ groups by average linkage \citep{monti2003consensus}. The exemplar of each group is the
member that affinity propagation chooses most often across the modal region, and the \emph{stability} is the mean
adjusted Rand index between those partitions and the consensus.

Finally, the distance can be explained. It is a sum of terms, one per direction and matched pair, so it splits over the
points of the two diagrams, and each point passes its share to the two seeds that realize its birth and its death: the
whole distance lies on critical points. The fields are closed-form functions of the trajectories, so each seed's share
is then divided among the agent positions in proportion to the size of the derivative of its value
\citep{leygonie2022framework,carriere2021optimizing}. The gradient of a field vanishes at its critical points inside
the domain, so a critical value does not change at first order when its point moves with the trajectories (the
envelope theorem), and holding the seeds fixed is exact at first order (\cref{app:method}). The method uses no label; to check what it finds, we apply it to robot swarms whose missions are known.

\FloatBarrier
% ============================================================================
\section{Experimental protocol}
\label{sec:protocol}

The runs are those of \citet{legarda2026universal}. In each run, $20$ e-puck robots \citep{mondada2009puck} are
simulated in ARGoS \citep{pinciroli2012argos}, in a dodecagonal arena of $4.91$\,m$^2$, and their positions are
recorded every $0.1$\,s. All the robots of a run execute the same controller, a probabilistic finite-state machine
designed by AutoMoDe-Chocolate \citep{francesca2014automode,francesca2015chocolate}. \citet{hasselmann2021empirical}
designed ten such controllers for each of five missions, and each mission asks for a different motion
(\cref{fig:missions}). In \emph{Directional Gate} (120\,s), the robots must traverse a gate at the center of the arena
from north to south, guided by a black corridor. In \emph{Foraging} (180\,s), they must reach one of two black food
sources and return to a white nest, with a light source as a cue. In \emph{Homing} (120\,s), they must gather on a
black area, starting from the upper part of the arena. In \emph{Shelter} (180\,s), they must aggregate in a white
shelter at the center, closed on three sides and lit from the open side; two black areas beside it play no role in
the mission. In \emph{XOR Aggregation} (180\,s), they must all aggregate on one of two black areas.

\begin{figure}[!htbp]
\centering
\figin{\linewidth}{figures/fig_missions.pdf}
\caption{The five missions seen from above (north to the left, light source in yellow), with the tracks of the
exemplar run of each mission.}
\label{fig:missions}
\end{figure}

Each controller was simulated in $50$ runs, with the initial conditions and durations of
\citet{hasselmann2021empirical}, and the ten controllers of each mission are ranked by their mean performance over
these runs. We call Top-$n$ the dataset that pools the runs of the $n$ best controllers of every mission, for $250n$
runs. Top-1, the $250$ runs of the best controller of each mission, is the dataset of that earlier study. Pooling
controllers adds behavioral diversity, since different controllers can solve a mission with different motion
strategies, and this diversity is what the method should discover.

The fields, diagrams and distances are computed once for the \nm{n-runs} runs and their \nm{n-pairs-million} million
pairs, and every dataset is drawn from them. Because the method is unsupervised and does not know how many controllers or
runs belong to each behavior, the taxonomy is read on a fixed local grid combining nine standard SNF neighborhood sizes
$\knn\in\{3,5,\dots,25\}$ ($1.2\%$ to $10\%$ of Top-1) with eleven diffusion lengths $t\in\{1,\dots,25\}$. To map the
full hierarchy, every dataset is also swept over a wide window proportional to its size $N=250n$: $\knn$ takes the
values $1, n, 2n, \dots, 50n$, reaching a fifth of the dataset ($N/5$), and $t$ runs from $1$ to $\nm{sw-t-max}$ (with
$\knn<3$ omitted in the state maps, where a run connects only to itself or one neighbor). Here $n$ only indexes the
dataset size, without passing labels or group counts to the clustering. No constant is tuned on the mission labels;
\cref{app:settings} lists them and varies each alone. Under the author's direction, a large language model was used to
review the claims in the paper, verify that the equations match the CPU and GPU implementations, and help accelerate
the GPU computation (see the declarations).

The mission labels are used only to score the taxonomies. With the $N$ runs of a dataset in clusters
$G_1,\dots,G_C$ and missions $L_1,\dots,L_5$, the purity $\frac{1}{N}\sum_{q}\max_{j}|G_q\cap L_j|$ is the fraction of
runs that lie in a cluster dominated by their own mission; the other runs are \emph{misassigned}. Since splitting alone
raises the purity, it is read with the number of clusters and with two scores based on entropies
\citep{rosenberg2007vmeasure}: the homogeneity $1-H(L\mid G)/H(L)$, which is one when no cluster mixes missions, and the
completeness $1-H(G\mid L)/H(G)$, which is one when no mission is split. The adjusted Rand index
\citep[ARI;][]{hubert1985comparing} compares the two partitions as a whole; it is one when they coincide and zero on
average by chance. Missions, however, are not the ground truth of the behaviors. A mission solved by several
controllers can contain several behaviors, whose discovery is the purpose of the method, and completeness and ARI
count them against it. On Top-1, where each mission has a single controller, every score is expected to reach one.


\FloatBarrier
% ============================================================================
\section{Results}
\label{sec:results}

The results open with the global phase map of stable states across resolutions (\cref{sec:res_states}), which explains
why we zoom into a local grid to count fine behaviors (\cref{sec:res_control}), and then measure how the coarse
transitions scale with dataset size (\cref{sec:res_law}). The exemplars, the attribution and the contribution of each
part of the method follow (\crefrange{sec:res_exemplars}{sec:res_why}).

\subsection{Stable states across resolutions}
\label{sec:res_states}

Over the wide proportional window, the resolution sweep maps the nested levels of the taxonomy (\cref{fig:states}).
Each map is made of a few broad stable states, and these states are nested. From small to large neighborhoods, the
runs pass from pieces of the missions to the five missions, then to Shelter merged with XOR Aggregation. Beyond lie
three families: Directional Gate with Foraging, Homing alone, and Shelter with XOR Aggregation. On Top-1, the five
missions hold \nm{sw-top1-missions} of the \nm{sw-n-settings-top1} settings and the three families
\nm{sw-top1-families}. There the states are stable in the strict sense: without any label, \nm{sw-top1-same-pct}\,\% of
the settings return exactly the partition of their neighbors.

On the pooled maps, the five missions hold fewer settings as controllers are pooled, from \nm{sw-missions-min-1to7} to
\nm{sw-missions-max-1to7} up to Top-7, then \nm{sw-missions-top8}, \nm{sw-missions-top9} and \nm{sw-missions-top10} on
Top-8, Top-9 and Top-10. Partial merges take their place: a block of runs, often most runs of one controller, joins
another mission before its own mission has gathered. Meanwhile the coarse three-family region expands to fill most of
the wide window (\nm{sw-families-min-pooled} to \nm{sw-families-max-pooled} settings on every pooled map). Selecting the
global mode across this full window would therefore collapse pooled datasets to the three coarse families, which is why
we zoom into the dotted local grid in \cref{sec:res_control}, where \nm{sw-maingrid-finer-min} to
\nm{sw-maingrid-finer-max} of the $99$ settings lie in the fine regime that preserves sub-behaviors.

\begin{figure}[!b]
\centering
\figin{\linewidth}{figures/fig_states.pdf}
\caption{Stable states across resolutions on each dataset over a window proportional to its size: $\knn$ from the first
step at least $3$ up to $50n$ and $t$ from $1$ to $\nm{sw-t-max}$. The color gives the state read from the clusters mixed
above a tenth (single-run clusters set aside, leaving four groups in Shelter + XOR and three or two in the coarse
family phases); hatched: a mission partly merged with another. Black lines mark transitions, the dotted box the local
grid of \cref{fig:resolution}, and the dashed line the edge of purity one.}
\label{fig:states}
\end{figure}

The transitions are announced by a few runs. On Top-1, the purity is exactly one at \nm{sw-top1-pure} settings of the
swept grid. Beyond them, and before two missions merge, at most \nm{sw-top1-cross-max} runs cross first: XOR
Aggregation runs join the cluster of Shelter, and Directional Gate runs that of Foraging, the two pairs that merge
next.

Homing never shares a cluster with another mission up to Top-\nm{sw-homing-free-max}. Because its robots first cross the
arena and then gather on a single patch, it combines the arena-wide kinetic intensity $K_\sigma$ of the traffic
missions with the concentrated density $\rho_\sigma$ and sink $\varphi_\sigma$ of the aggregation missions, keeping it
apart from both families in the distance space. On Top-9 and Top-10, it merges only beyond
\nm{sw-homing-knn-pct}\,\% of the largest neighborhood, at \nm{sw-two-top9} and \nm{sw-two-top10} settings. There it
joins Directional Gate and Foraging in a last state of two families, apart from Shelter and XOR Aggregation.

The pairs that merge are not those of equal duration, since Directional Gate (120\,s) joins Foraging (180\,s) and not Homing
(120\,s). They are the missions that move alike, and the two pairs that no field separates well (\cref{fig:pairs}).
Directional Gate and Foraging are both traffic across the arena, distinguished mainly by how this traffic turns, which
the vorticity potential separates best. Shelter and XOR Aggregation both end with the robots gathered on a patch,
distinguished mainly by where the flux converges, the domain of the divergence potential. Some XOR Aggregation runs even
keep moving until near their end, as Shelter runs do (\cref{app:merge}). Homing, which stays apart, differs from each
family in another field: from the traffic mainly in its divergence potential, and from the patches mainly in its
kinetic intensity. The coarse states of the map thus group behaviors by the organization of their motion.

\begin{figure}[!htb]
\centering
\figin{\linewidth}{figures/fig_pair_separation.pdf}
\caption{Separation of each pair of missions by each field on Top-1, the distance between missions over the distance
within them; bold: the best field of each pair.}
\label{fig:pairs}
\end{figure}

\subsection{Zooming into the local scale: how many behaviors}
\label{sec:res_control}

Without being told how many behaviors there are, zooming into the local grid recovers the five missions of Top-1
(\cref{fig:resolution}). Five clusters hold over most of the local grid, \nm{top1-plateau} of the $99$ settings, the
largest modal region of the ten datasets, and \nm{top1-errfree} of the $99$ settings misassign no run. Off the modal
region, only the finest settings give more clusters, and these cut the missions into pure pieces rather than mixing
them.

Pooling controllers adds sub-behaviors rather than errors. Across the ten local maps, purity stays at or above $0.995$
at \nm{all-pur-ge995} of all \nm{all-settings-n} settings (\nm{all-pur-ge995-pct}\,\%): a mission solved by several
controllers divides into controller-level sub-behaviors instead of mixing with another mission. On Top-10,
\nm{top10-err-le2} of the $99$ settings misassign at most two of the $2{,}500$ runs. Only the small-neighborhood
transient bridge at $\knn=\nm{top10-worst-knn}$ for $t=5$--$9$ drops to purity \nm{top10-worst-pur}
(\nm{top10-worst-err} misassigned runs at $t=\nm{top10-worst-t}$), before longer diffusion restores purity at or above
\nm{top10-k3-t10plus-pmin} for all $t\ge 10$.

\begin{figure}[p]
\centering
\figin{\linewidth}{figures/fig_resolution_map.pdf}
\caption{Local resolution maps of the ten datasets over $\knn\in[3,25]$ and $t\in[1,25]$: number of clusters $C$ at each
setting, purity against the five missions on $[0,1]$ (at least \nm{maps-purity-min} across all \nm{all-settings-n}
settings), and the modal region ($C=\Cstar$) outlined in white.}
\label{fig:resolution}
\end{figure}
% the maps fill a page of their own: the next paragraph, on the table, starts after them
\clearpage

Scored against the missions (\cref{tab:taxonomy}), the modal consensus of Top-1 places every run with its mission. The
taxonomies of the pooled datasets keep a homogeneity close to one, since no cluster mixes missions, while their
completeness and ARI decrease as controllers are pooled. They decrease for this reason alone, since both scores count
the discovered sub-behaviors against the method: once the clusters that share a dominant mission are merged, the
missions come back almost exactly. The few runs misassigned over all datasets, \nm{all-misassigned} of \nm{all-runs},
fall between the two pairs of missions that merge first at coarse resolution, a first sign that they reflect a
resemblance of motion rather than noise. Similarly, Top-8 to Top-10 have no zero-error entry in \cref{tab:taxonomy} only
because at least one boundary run out of $2{,}000$ to $2{,}500$ crosses into the mission whose motion it shares at every
setting of the local grid.

\begin{table}[!htb]
\centering
\caption{Taxonomy of each dataset from the local grid: without labels, the modal count $\Cstar$, its modal region and
its stability; against the missions, the scores of the modal consensus and the settings without any misassigned run
(those in the modal region in parentheses).}
\label{tab:taxonomy}
\tabin{taxonomy}
\end{table}

The sub-behaviors are not spread evenly over the missions. XOR Aggregation divides first, from Top-2, and always into at
least as many clusters as any other mission, while the others divide only as more controllers are pooled. Nor do the
sub-behaviors follow the controllers one to one. A sub-behavior can be exactly the runs of one controller, as for
\nm{iso-mission} on Top-\nm{iso-top}, but on Top-10 the runs of \nm{top10-ctrl-split} of the \nm{top10-ctrl}
controllers fall into more than one cluster, and most clusters gather several controllers of a mission. The number of
behaviors stops growing at \nm{top10-cstar}, from Top-8 to Top-10, far below the number of controllers pooled.

\subsection{The scaling of the transitions}
\label{sec:res_law}

The transitions move with the size of the dataset, since a neighborhood of fixed size covers a smaller part of a larger
mission. Measured in units of $n$ (with $N=250n$ trajectories), as in the window of \cref{fig:states}, they align across
the regular pooled datasets. Across all six regular pooled datasets (Top-2 to Top-7), the transition to the three
families follows one power law, $\knn/n=\nm{law-c}\,t^{-\nm{law-alpha}}$ (\cref{fig:law}), and the exponent fitted on
each dataset alone stays between \nm{law-fit-alpha-min} and \nm{law-fit-alpha-max}. The exponent is close to the one half
of a random walk. The reach of the diffusion would then grow like $\knn\sqrt{t}$, and the families would merge once it
covers a fixed fraction of a mission.

The scaling has bounds. Top-1, whose missions are single controllers, follows a shallower power law,
$\knn/n=\nm{law-top1-c}\,t^{-\nm{law-top1-alpha}}$. From Top-8 on, the exponent steepens, to between
\nm{law-drift-alpha-min} and \nm{law-drift-alpha-max}, as the datasets take in the weakest controllers, some of which do
not move like their mission. On Top-10, the runs of the tenth Shelter controller sit with Directional Gate or Foraging
at \nm{drift-s10-nav} of the \nm{sw-n-settings} settings. Their robots follow a strategy atypical of their mission: they keep
moving, at a median speed of
\nm{drift-s10-speed}\,m/s, against \nm{drift-shelter-speed-min} to \nm{drift-shelter-speed-max}\,m/s for the nine other
Shelter controllers and \nm{drift-nav-speed}\,m/s for the traffic missions. The runs of the tenth XOR Aggregation
controller, for their part, sit with Shelter at \nm{drift-x10-shelter} settings. The map places these runs with what
their robots do, not with the mission they were given, and so singles out, without labels, the controllers whose robots
do not behave like those of their mission.

\begin{figure}[!htb]
\centering
\figin{\linewidth}{figures/fig_law.pdf}
\caption{Scaling of the transition to the three families. (a)~The value of $\knn/n$ at which each dataset first reads
the three families at each $t\geq5$, on log-log axes; dashed: the power law fitted on Top-2 to Top-7; dotted: the power
law fitted on Top-1. (b)~Exponent fitted on each dataset alone; shaded: range of Top-2 to Top-7; open orange circles:
Top-8 to Top-10, which include atypical controllers.}
\label{fig:law}
\end{figure}

\subsection{One exemplar per behavior}
\label{sec:res_exemplars}

The map says how the behaviors group at every resolution, but not what a behavior looks like. For this, the taxonomy gives each behavior an exemplar, the run that affinity propagation selects most often
across the modal region (\cref{fig:fields}). An exemplar is a recorded run, not an average of runs, so it can be replayed,
inspected or imitated, and its fields show the behavior directly. Each mission leaves its own imprint on them. In
Directional Gate, the density and the kinetic intensity follow the corridor and the gate, and the divergence potential
splits the arena into a source half and a sink half. Foraging fills most of the arena with round trips. Homing combines
a single dense spot, where the robots end, with an agitation spread over the arena by the crossing that precedes it.
Shelter shows the shelter and the lit wall in its density, and bands of opposite rotation on either side of the
shelter. XOR Aggregation, finally, splits its robots between its two black areas.

\begin{figure}[!htb]
\centering
\figin{0.8\linewidth}{figures/fig_fields.pdf}
\caption{The five fields of the exemplar of each behavior on Top-1, at $\sigma_3=\nm{sigma3-cm}$\,cm, beside the
tracks of its robots (final positions in black). Density and kinetic intensity run from light to dark blue, the
potentials and the coupling from blue (negative) through grey to red (positive), each panel on its own scale. A red
maximum of $\varphi_\sigma$ is a sink and a blue minimum a source; a blue minimum of $\psi_\sigma$ is a
counterclockwise rotation center. White: farther than $5\sigma$ from every robot.}
\label{fig:fields}
\end{figure}

The five fields see different things. The density and the kinetic intensity coincide where the robots keep moving, as
in Directional Gate and Foraging, and part where they stop: the density of Homing marks where its robots end, while its
kinetic intensity keeps the crossing that brought them there. A robot at rest adds to the density only, since the
kinetic intensity, the two potentials and the coupling all vanish with its velocity. The potentials, for their part,
show the direction of the motion, toward sinks and around rotation centers, which the density and the kinetic intensity
ignore. Distinct as they are, some of these imprints resemble one another, as the map showed, and the attribution of the
distance says in what.

\subsection{Why two behaviors differ}
\label{sec:res_attribution}

The last component of the taxonomy explains its differences. Splitting the distance between two exemplars, term by
term over the diagram points and from each point over the seeds of its birth and death, shows what makes them differ. We follow Directional Gate and Foraging as an example (\cref{fig:attribution}), before the ten pairs of
missions (\cref{tab:attribution}). On the vorticity
potential, which separates them best, the whole distance lies on critical points, and the two rotation centers of the
traffic carry the largest share, \nm{attr-centres-exemplar}. The rest spreads over the other critical points, saddles
included, so the difference rests on the organization of the field rather than on a single structure. Split through the
distributions of values instead, \nm{attr-val-ordinary} of the distance lands on ordinary points that merely hold a
quantile, and the split is less stable under noise. Carried back to the robots, the difference concentrates at the end
of the corridor, where the robots gather and turn. Repeated over the \nm{attr-runs} runs of the Directional Gate
controller, the largest shares fall on the same two rotation centers in every run, and their mean map keeps two compact
regions, whereas the split of the values spreads over the whole positive cell of the potential.

\begin{figure}[!htb]
\centering
\figin{0.78\linewidth}{figures/fig_attribution.pdf}
\caption{The distance between the Directional Gate and Foraging exemplars split over the seeds, carried back to the
robots and averaged over \nm{attr-runs} runs, for extended persistence (a--c) and the distribution of values (d--f).
Single-run robot maps (b, e) use square-root color scaling with robot trajectories overlaid in light gray.}
\label{fig:attribution}
\end{figure}

The example is not an isolated case. Split on the field that separates each pair best, every pair of missions places
its whole topological distance on critical points (\cref{tab:attribution}). Across the \nm{attr-runs} runs of the first
controller, \nm{attrp-pairs-50} of the ten pairs place their two largest extrema on a maximum and a minimum in every
single run ($50/50$). Only Directional Gate against XOR Aggregation drops to $\nm{attrp-recur-min}/50$, because a single
run splits its largest weights across two maxima of the density $\rho_\sigma$ rather than pairing the maximum with a
minimum. These two extrema carry a larger share of the split on the potentials and their coupling than on the density
and the kinetic intensity, whose differences spread over more critical points. Under a noise of $2$\,mm, the topological
split stays on the same seeds more than that of the values for every pair. Two mechanisms explain why the exact-seed
score drops to $0.65$--$0.66$ on three pairs: on Homing paired with Shelter or XOR Aggregation via kinetic intensity
$K_\sigma$, $2$\,mm jitter on nearly resting robots changes the small peak value of $K_\sigma$ by up to
\nm{attrp-homing-k-change}; on Shelter against XOR Aggregation via $\varphi_\sigma$, critical seeds shift locally
between neighboring vertices while \nm{attrp-sx-blur} of the topological weight stays within one footprint width
$\sigma$. A difference between two behaviors is thus named by the structures that carry it and located on the robots
that shaped them.

\begin{table}[!htb]
\centering
\caption{The split of the distance between the Top-1 exemplars of each pair on its best field (\cref{fig:pairs}):
share of the two largest extrema, runs where they recur, share kept under $2$\,mm of noise.}
\label{tab:attribution}
\tabin{attribution_pairs}
\end{table}

\subsection{What each part of the method contributes}
\label{sec:res_why}\label{sec:res_compare}\label{sec:res_robust}

The four components rest on three parts of the method: the fields, the reading of several scales together and the
topology of the fields. Removing each in turn shows what it contributes. No field separates the missions without error
on its own: each alone misassigns runs at every one of the $990$ settings of the ten datasets, and the five together
reach the highest purity and the most settings without any misassigned run (\cref{fig:ablation}). The density is the
only field to which a robot at rest contributes, the other four vanishing with its velocity, and without it the settings
without error fall from \nm{abl-all-errfree} to \nm{abl-KphipsiJ-errfree}. The coupling adds what the separate diagrams
of the two potentials miss, where the flux converges relative to where it turns: with it, Top-1 has
\nm{abl-top1-all-errfree} settings without any misassigned run instead of \nm{abl-top1-noJ-errfree}, and the number of
behaviors never decreases when controllers are added. Their median ARI notches in \cref{fig:ablation} overlap
(\nm{abl-all-ari-median} against \nm{abl-noJ-ari-median}), because without the coupling some pooled settings merge
sub-behaviors into fewer groups against the five missions.

\begin{figure}[!t]
\centering
\figin{\linewidth}{figures/fig_ablation.pdf}
\caption{Field ablation over the $99$ settings of the ten datasets: purity and ARI of each combination of fields, with
the median and, for purity, the number of the $990$ settings without any misassigned run.}
\label{fig:ablation}
\end{figure}

To isolate what the representation contributes, other descriptions of the runs are clustered with the same fusion,
affinity propagation and modal consensus. One summarizes each run by the means and the deviations of four order
parameters: the polarization, the milling, the mean speed and the distance to the centroid. Another uses
the Betti curves of the robot positions, as CROCKER plots do \citep{topaz2015topological,bhaskar2019analyzing}. The
last keeps our fields and seeds but describes each field by the distribution of its values instead of its diagrams,
which isolates what the topology adds. Two distributions are compared by the mean absolute difference of
$\nm{n-quantiles}$ evenly spaced quantiles, an estimate of the Wasserstein distance of order one on the line, so that
the two summaries are compared at the same order. Top-1 is also the dataset that the earlier study grouped into five
clusters. From the confusion matrix it reports, with $50$ runs per mission, we recover the counts of that partition and
score it as ours.

The descriptions built without the fields fall short (\cref{tab:baselines}, and
\cref{tab:baselines_full} for all the runs). Order parameters, which average the motion over the whole group, mix
Directional Gate with Foraging and split XOR Aggregation, and CROCKER curves, which see where the robots are but not how
they move, place all Foraging runs with Directional Gate. Neither leaves a single one of the $99$
settings without a misassigned run, where the fields leave \nm{base-ours-top1-free}. Reading the scales together matters as well: without
diffusion, the fields split Top-1 into \nm{base-single-top1-cmin} to \nm{base-single-top1-cmax} nearly pure groups,
which the fusion across scales joins into the five missions. The topology, finally, sharpens the separation on diverse controllers and brings the explanation. On Top-1, the
distributions of the field values separate the five missions almost as well as the diagrams (\cref{tab:baselines}). On
Top-10, however, the diagrams retain higher purity (\nm{base-ours-top10-purity} against \nm{base-fieldstat-top10-purity}),
homogeneity (\nm{base-ours-top10-hom} against \nm{base-fieldstat-top10-hom}), completeness (\nm{base-ours-top10-comp}
against \nm{base-fieldstat-top10-comp}) and ARI (\nm{base-ours-top10-ari} against \nm{base-fieldstat-top10-ari};
\cref{tab:baselines_full}). Combining diagrams with values yields \nm{base-combined-top10-free} error-free settings out
of $99$. Above all, only the diagrams split the distance over the critical structures and the robots that carry it
(\cref{sec:res_attribution}).

\begin{table}[!t]
\centering
\caption{Descriptions of the $250$ runs of Top-1 under the same fusion, affinity propagation and modal consensus (Top-10
in \cref{tab:baselines_full}); the method is shaded. $^{*}$Partition published for the same runs, with five fixed
clusters \citep{legarda2026universal}.}
\label{tab:baselines}
\tabin{baselines_top1}
\end{table}

The taxonomy is also stable. It does not depend on what it should ignore: rotating the scene changes the diagrams far
less than two runs of the same mission differ (\cref{app:invariance}). Nor does it depend on the details of its
construction: varied alone, no constant of the method changes the number of behaviors on Top-1 (\cref{app:settings}),
and each step of the fusion is needed, since without the mutual ranks, for instance, XOR Aggregation runs move to
Shelter (\cref{app:fusion}).

% ============================================================================
\section{Discussion, limitations and future work}
\label{sec:discussion}\label{sec:limits}\label{sec:future}

The results separate three notions that are often taken as one. A mission is defined by the objective that scores it,
a controller by the rules that drive its robots, and a behavior, here, by the organization of the motion. Here they
coincide only on Top-1, where each mission has one controller. Once controllers are pooled, a mission divides into several
behaviors, a controller can produce more than one behavior across its runs, and several controllers can share one, so
that the \nm{top10-ctrl} controllers of Top-10 give \nm{top10-cstar} behaviors. At coarse resolution, in turn, behaviors
of different missions join into families of motion: traffic across the arena, gatherings on patches, and Homing apart.
Scores against the missions count these discoveries as errors, which is why completeness and ARI decrease while purity
does not. Missions and controllers are two labelings of the motion, and each describes it only in part.

% this paragraph is never split across two pages: it moves whole to the next page when it does not fit
{\interlinepenalty=10000 The parts of the method also play distinct roles. The fields and the reading of several
scales together separate the behaviors. On the same fields, a distribution of values reaches almost the same purity
and homogeneity on Top-10 (\nm{base-fieldstat-top10-purity} and \nm{base-fieldstat-top10-hom}, against
\nm{base-ours-top10-purity} and \nm{base-ours-top10-hom} for the diagrams). The diagrams, however, retain higher
completeness and ARI (\nm{base-ours-top10-comp} and \nm{base-ours-top10-ari}, against \nm{base-fieldstat-top10-comp}
and \nm{base-fieldstat-top10-ari}), and they both separate and explain the behaviors. Their distance splits over the
critical points that carry it, such as sinks, sources and rotation centers of the potentials. This split stays on the
same structures under noise and from run to run. Since the gradient of a field vanishes at its interior critical points,
the split also carries back to the robots. A description that only separated the behaviors would answer half of the
question of the introduction: which behaviors there are, but not what distinguishes them.\par}

This pairing of an exemplar with an explainable distance is what an imitation needs. A grouping of recordings says
which behaviors exist, but not what to change in a controller to reproduce one of them. Here the exemplar is a recorded
run that a controller can be asked to reproduce. The distance to it scores an attempt. When the attempt fails, it also
names the missing structure and where in the motion it lies.

The measure does not see everything in the motion, and what it misses bounds what the results show. The most
important limit is the average over the recording, which erases the chronology. The XOR Aggregation runs that resemble
Shelter keep moving for most of the run and come to rest only near its end (\cref{app:merge}). An average over the
whole run cannot see when a behavior settles.

A second limit is that the method observes trajectories only. It cannot distinguish two missions that produce the same
motion in response to different stimuli. Aggregating on a black patch and aggregating on a white patch, for example,
can produce the same motion, because neither the sensors of the robots nor the patch color enters the fields.

These limits point to the next steps of this work. The first is to apply the same method to the sensors of the robots.
The footprint, the decomposition, the coupling and extended persistence apply without modification to sensor
measurements projected in space. Imitation comes after. Our previous work imitated behaviors with a topological
similarity of motion fields \citep{ndinga2025thesis}. For transfer between two controller architectures, it still
combined this similarity with mission performance. Combining the topological objective on the motion with the same
objective on the sensors makes another aim possible. A controller can then imitate a behavior to accomplish a mission
without any definition of task performance. Extending the representation in time is a natural next step, and with time the method can
reach beyond robotics. A description that follows a collective in time could detect automatically the transitions of a
biological process, such as the division of cells by mitosis, and point to what changes during them. The footprint
itself can be enriched: beyond a round Gaussian, it can take the shape of the agents it represents, such as real cells,
while keeping the whole computation differentiable.

\FloatBarrier
% ============================================================================
\section{Conclusion}
\label{sec:conclusion}

We asked which behaviors a set of recordings of a collective contains, and what in the motion distinguishes them. We
answered with a taxonomy recovered from the motion alone, without labels: the number of behaviors, the recordings that
share each, one exemplar per behavior, their hierarchy across resolutions and the attribution of their differences.
Five rotation-invariant fields, summarized by extended persistence without a grid at scales derived from the data, are
compared by sliced Wasserstein distances, fused across scales and clustered by affinity propagation. On a robot swarm,
the method separates the five missions without a misassigned run when each has one controller, reveals the strategies
of pooled controllers as pure sub-behaviors, and merges at coarse resolution the missions that are close in motion. The
fields separate the behaviors, the multiscale view read together by the fusion gathers the recordings of each, and the
topology provides the measure: a distance between organizations of the motion that differentiation traces back to the
robots. It is the measure needed to test whether an imitation reproduces a behavior.

\FloatBarrier
\section*{Declarations}

\bmhead{Funding} The author declares that no funds, grants or other support were received during the preparation of
this manuscript.

\bmhead{Competing interests} The author has no relevant financial or non-financial interests to disclose.

\bmhead{Ethics approval and consent to participate} Not applicable: the study uses simulated robot data only.

\bmhead{Consent for publication} Not applicable.

\bmhead{Data availability} The runs analyzed here are those of \citet{legarda2026universal}, whose trajectories are
publicly available at \url{https://github.com/glegarda/swarm-tda}. The distance matrices, the clustering results and
every number reported here are deposited on Zenodo (\doilink{10.5281/zenodo.23440748}).

\bmhead{Code availability} The code that reproduces every result and figure, from the trajectories to this article,
is available at \url{https://github.com/Vertu5/paper_topological_kinematics} and archived on Zenodo
(\doilink{10.5281/zenodo.23440748}).

\bmhead{Use of large language models} Under the author's direction, Claude (Anthropic) was used to review the claims in
the manuscript, to verify that the equations in the paper match both the CPU and GPU implementations, and to help
accelerate the GPU computation. Every check and result was verified by the author. The research question, the method,
the design of the experiments, the interpretation of the results and all scientific claims are the author's, as is full
responsibility for the content.

\bmhead{Author contributions} Olivier Vertu Ndinga Oba designed the study, developed the method, wrote the first
implementation of the code and the first draft of the manuscript, designed the analyses and the plots, interpreted the
results, and verified and approved the final code and text (conceptualization, methodology, software, validation,
formal analysis, investigation, data curation, visualization, writing: original draft, review and editing).

\ifdefined\journalversion\else
\clearpage
\begin{appendices}
% The appendices, in the order in which the paper first cites them. paper.tex (arXiv) inputs them inside its
% appendices environment; supplementary.tex inputs them as the Supplementary Information of the journal version.
% ============================================================================
\section{Proofs and verification}
\label{app:proofs}

This appendix proves that the closed forms \eqref{eq:potentials} solve the Helmholtz--Hodge decomposition, proves the
two other propositions and gives the closed forms that the code evaluates.

\begin{proposition}[Closed-form Helmholtz--Hodge potentials]\label{prop:hhd}
For a single agent, $\nabla\big(u\,\kappa_\sigma(\vd)\big)-\nabla^{\perp}\big(w\,\kappa_\sigma(\vd)\big)
=\vv\,e^{-\|\vd\|^2/2\sigma^2}$ at every $\vd\neq\vec 0$. Consequently the potentials \eqref{eq:potentials} satisfy,
at every point $\vr$,
\[
\nabla\varphi_\sigma+\nabla^{\perp}\psi_\sigma=\frac{1}{NT\sigma\vbar}\sum_{t}\vV_\sigma(\vr,t).
\]
\end{proposition}

\begin{proof}[Proof of \cref{prop:hhd}]
The idea is to invert the Laplacian on the flux of one agent, that is, to solve the Poisson equation of the plane
for a vector field. The solution is then split into a gradient and a rotated gradient, and the two potentials are
the scalar fields that appear in this split.

Consider one agent at the origin, with a constant velocity $\vv$, and write $\varrho=\|\vd\|$ for the
distance to it. Its contribution to the flux is $\vv\,e^{-\varrho^2/2\sigma^2}$, a Gaussian bump carried in
the direction of $\vv$. We first look for a radial function $\Gamma(\varrho)$ whose Laplacian is this
Gaussian, up to a factor $1/\sigma^2$. For a radial function, the Laplacian reads
$\Delta\Gamma=\varrho^{-1}(\varrho\Gamma')'$. Integrating from the origin, where $\Gamma$ must stay smooth, gives
\[
\varrho\,\Gamma'(\varrho)=\int_0^{\varrho}\frac{s}{\sigma^2}\,e^{-s^2/2\sigma^2}\,\mathrm{d}s=1-e^{-\varrho^2/2\sigma^2},
\]
which gives
\[
\Delta\Gamma=\frac{1}{\varrho}\,\frac{\mathrm{d}}{\mathrm{d}\varrho}\Big(1-e^{-\varrho^2/2\sigma^2}\Big)
=\frac{1}{\sigma^2}\,e^{-\varrho^2/2\sigma^2}.
\]
Far from the agent, $\Gamma'\simeq1/\varrho$, so $\Gamma$ grows like the logarithm of the distance, as the free-space
Green's function of the plane does \citep{bhatia2014natural}. Near the agent, the Gaussian smooths it.
Compared with \eqref{eq:kappa}, this choice reads $\sigma^2\Gamma'(\varrho)=\varrho\,\kappa_\sigma(\varrho)$:
the kernel $\kappa_\sigma$ is $\sigma^2$ times the radial derivative of $\Gamma$, divided by the
distance.

Second, every smooth vector field $\vW=(W_x,W_y)$ of the plane satisfies
\[
\Delta\vW=\nabla(\nabla\!\cdot\vW)+\nabla^{\perp}(\nabla^{\perp}\!\cdot\vW),\qquad
\nabla^{\perp}\!\cdot\vW=\partial_xW_y-\partial_yW_x,
\]
as the expansion of both sides component by component shows. We apply this identity to
$\vW=\sigma^2\Gamma\,\vv$. Since $\vv$ is constant, $\Delta\vW=\sigma^2(\Delta\Gamma)\,\vv=\vv\,e^{-\varrho^2/2\sigma^2}$,
which is the contribution of the agent. The identity splits it into the gradient of
$\Phi=\nabla\!\cdot\vW=\sigma^2\,\vv\cdot\nabla\Gamma$ and the rotated gradient of
$\Psi=\nabla^{\perp}\!\cdot\vW=\sigma^2\big(v_y\,\partial_x\Gamma-v_x\,\partial_y\Gamma\big)$:
\[
\vv\,e^{-\varrho^2/2\sigma^2}=\nabla\Phi+\nabla^{\perp}\Psi .
\]

Third, we write $\Phi$ and $\Psi$ with the kernel. The gradient of a radial function points along $\vd$,
and $\sigma^2\nabla\Gamma=\sigma^2\Gamma'(\varrho)\,\vd/\varrho=\kappa_\sigma(\vd)\,\vd$. Hence
$\Phi=\kappa_\sigma(\vd)\,(\vv\cdot\vd)=u\,\kappa_\sigma(\vd)$ and
$\Psi=\kappa_\sigma(\vd)\,(v_yd_x-v_xd_y)=-w\,\kappa_\sigma(\vd)$, which is the single-agent identity of the
proposition. The computation holds wherever $\vd\neq\vec 0$. At the agent itself both sides are smooth, and
the identity follows by continuity.

Finally, the identity is linear in the velocity. Summing it over the agents and the instants, with the
factor $1/(NT\sigma\vbar)$, gives the potentials \eqref{eq:potentials} on the left and
$\frac{1}{NT\sigma\vbar}\sum_t\vV_\sigma(\vr,t)$ on the right, since $\vV_\sigma=\sum_i\vv_iO_i$. The two
potentials therefore rebuild the time-averaged flux exactly, without a grid and without a boundary
condition.
\end{proof}

\begin{proposition}[Invariance]\label{prop:so2}
Let a rigid transformation of the plane, a rotation $Q\in\mathrm{SO}(2)$ followed by a translation, act
on all positions, on all velocities and on the evaluation point. Then $\rho_\sigma$, $K_\sigma$, $\varphi_\sigma$,
$\psi_\sigma$ and $J_\sigma$ are unchanged, while $\nabla\varphi_\sigma$ and $\nabla\psi_\sigma$
are multiplied by $Q$.
\end{proposition}
\begin{proof}
Every field is built from a few elementary quantities: the distances $\|\vd_i\|$ between the point and
the agents, the speeds $\|\vv_i\|$, the scalar products $u_i=\vv_i\cdot\vd_i$, the determinants
$w_i=\det[\vv_i,\vd_i]$, and the averages $\vbar$ and $\overline{v^2}$. A translation changes neither $\vd_i$
nor $\vv_i$. A common rotation preserves lengths, angles and the orientation of the plane, so it changes
none of these quantities either. Hence $\rho_\sigma$, $K_\sigma$, $\varphi_\sigma$ and $\psi_\sigma$ take
the same value at the transformed point, and their gradients turn with the scene: they are multiplied
by $Q$. The coupling keeps its value, since a rotation does not change a determinant:
$\det[Q\nabla\varphi_\sigma,Q\nabla\psi_\sigma]=\det Q\,\det[\nabla\varphi_\sigma,\nabla\psi_\sigma]
=\det[\nabla\varphi_\sigma,\nabla\psi_\sigma]$, because $\det Q=1$.
\end{proof}
The seeds are selected from the agents and from the fields themselves, so the diagrams inherit this
invariance, up to floating-point effects in the choice of the seeds, which \cref{app:invariance}
measures.

\begin{proposition}[Tail of the footprint]\label{prop:tail}
For $O_i(\vr)=\exp(-\|\vr-\vx_i\|^2/2\sigma^2)$ and any $\varrho\ge0$,
$\int_{\|\vr-\vx_i\|>\varrho}O_i(\vr)\,\mathrm{d}\vr=2\pi\sigma^2e^{-\varrho^2/2\sigma^2}$. At $\varrho=5\sigma$
the fraction of the mass outside the evaluation window is $e^{-12.5}\approx3.7\times10^{-6}$.
\end{proposition}
\begin{proof}
The footprint depends only on the distance $s=\|\vr-\vx_i\|$ to the agent, so we integrate in polar
coordinates centered on it, where the element of area is $2\pi s\,\mathrm{d}s$:
\[
\int_{\|\vr-\vx_i\|>\varrho}O_i(\vr)\,\mathrm{d}\vr=\int_\varrho^\infty e^{-s^2/2\sigma^2}\,2\pi s\,\mathrm{d}s
=2\pi\sigma^2\Big[-e^{-s^2/2\sigma^2}\Big]_{\varrho}^{\infty}=2\pi\sigma^2e^{-\varrho^2/2\sigma^2}.
\]
With $\varrho=0$ the whole mass is $2\pi\sigma^2$, so the fraction left outside a window of radius
$\varrho$ is $e^{-\varrho^2/2\sigma^2}$. At $\varrho=5\sigma$ it is $e^{-25/2}=e^{-12.5}\approx3.7\times10^{-6}$.
\end{proof}

The coupling field needs the gradients of the two potentials. They follow from the same closed forms,
by differentiating $u_i\kappa_\sigma(\vd_i)$ and $w_i\kappa_\sigma(\vd_i)$ with respect to the point. With
$c=1/(NT\sigma\vbar)$, $\vv_i^{\perp}=(-v_{i,y},v_{i,x})$ and
$\lambda_i=\big(e^{-\|\vd_i\|^2/2\sigma^2}-2\kappa_\sigma(\vd_i)\big)/\|\vd_i\|^2$,
\[
\nabla\varphi_\sigma=c\sum_{t,i}\big(\kappa_\sigma(\vd_i)\,\vv_i+\lambda_i\,u_i\,\vd_i\big),\qquad
\nabla\psi_\sigma=-c\sum_{t,i}\big(\kappa_\sigma(\vd_i)\,\vv_i^{\perp}+\lambda_i\,w_i\,\vd_i\big).
\]
A seed taken from the agent positions coincides with an agent at some instant, so these expressions must remain
accurate at $\vd_i=\vec 0$. Two details ensure it. The code computes $1-e^{-s}$ as $-\operatorname{expm1}(-s)$,
which keeps the relative precision of $\kappa_\sigma$ as $\|\vd\|\to0$. It also floors $\|\vd\|^2$ at
$10^{-12}$\,m$^2$, where $\kappa_\sigma$ equals $1/2$ to better than the single precision in which the
fields are stored. The factor $\lambda_i$ tends to $-1/(4\sigma^2)$ as $\vd_i\to\vec 0$, but its numerator is
the difference of two numbers close to one, which loses precision. It is multiplied by $u_i\,\vd_i$ or
$w_i\,\vd_i$, of order $\|\vd_i\|^2$, so this loss does not reach the gradients. The derivatives with respect
to the positions (\cref{sec:res_attribution}) are taken by automatic differentiation of the same
expressions.

Every statement of this appendix is also checked by computer. The script
\nolinkurl{manuscript/verify_propositions.py} verifies with SymPy the single-agent identity of
\cref{prop:hhd}, the orthogonality relations $\nabla\cdot\nabla^{\perp}\psi=0$ and
$\nabla^{\perp}\cdot\nabla\varphi=0$, the invariance of the elementary terms of \cref{prop:so2}, and the
tail of \cref{prop:tail}, which a numerical quadrature confirms to $\nm{tail-quad-err}$. It also checks
the gradients above, the limit of $\lambda_i$ and their value at an agent, and the rotation invariance of
the magnitude of the traceless part of $\tM_2$ (\cref{sec:representation}). On a synthetic milling group, it
checks numerically the envelope property of \cref{app:attribution}.

% ============================================================================
\section{The method in detail}
\label{app:method}

This appendix gives the details of the construction of \cref{sec:method} and illustrates its steps on synthetic swarms;
\cref{fig:overview} follows the same steps on the robot data.

\subsection{The footprint and the moments}

The distribution of the agents is a set of points, which the footprint turns into continuous fields. Four properties
are asked of it. It must favor no direction, so that the fields do not depend on the orientation of the scene. It must
be smooth, so that the fields can be differentiated with respect to the positions of the agents. It must be negligible
beyond a finite window, so that a field is evaluated locally. And the potentials of the flux must follow from it in
closed form, by solving a Poisson equation without a grid. A Gaussian footprint has all four
(\cref{prop:hhd,prop:so2,prop:tail}). The same philosophy, observing a signal at all scales at once through Gaussians of
growing width, underlies the scale space of image analysis \citep{witkin1983scale,lindeberg1994scale}; here, however,
the Gaussian was chosen for the four properties above, and its width as a lens is only one of its roles. Up to a
constant factor, $M_0$ and $\vM_1$ are the density and the momentum density by which discrete particle systems are
coarse-grained into continua, where a Gaussian is also chosen for its smooth fields and exact integrals
\citep{weinhart2012discrete}.

Two footprints whose centers are $d$ apart are no longer resolved, in the sense that their sum at the midpoint exceeds
the height of a single footprint, when $d\le2\sqrt{2\ln2}\,\sigma\approx2.35\,\sigma$ (\cref{fig:footprint}). At the
smallest width of the ladder, $d_{\min}/2$, two agents are therefore resolved as soon as they are more than
$\sqrt{2\ln2}\,d_{\min}\approx1.18\,d_{\min}$ apart. The floor $\epsilon_\sigma=0.01$\,m of \eqref{eq:ladder} is the
tracking precision.

\begin{figure}[!htbp]
\centering
\figin{\linewidth}{figures/fig_footprint.pdf}
\caption{The Gaussian footprint: (a)~one agent, evaluated within $5\sigma$; (b)~its profile; (c)~two agents at
distance $d$, no longer resolved below $2.35\,\sigma$; (d,\,e)~the width as a lens, from the agents to the group.}
\label{fig:footprint}
\end{figure}
\FloatBarrier

Each moment captures something that the others miss (\cref{fig:moments}). A group at rest has $\vM_1=\vec 0$ but
occupies space in $M_0$. Two streams crossing in opposite directions cancel in $\vM_1$ but add up in
$\operatorname{tr}\tM_2=\sum_i\|\vv_i\|^2O_i$, the local kinetic intensity. Likewise, a concentrated group and a
dispersed one can share the same net flux and still differ in $M_0$. The traceless part of $\tM_2$, which measures how
strongly the velocities align along a common axis, is also invariant under rotations, but we leave it aside.

\begin{figure}[!htbp]
\centering
\figin{\linewidth}{figures/fig_moments.pdf}
\caption{The three moments on a synthetic configuration at one instant ($\sigma=6.5$\,cm): a group at rest, two
crossing streams and a group moving together (a), and $M_0$, $\|\vM_1\|$ and $\operatorname{tr}\tM_2$ (b--d).}
\label{fig:moments}
\end{figure}
\FloatBarrier

\subsection{The two potentials}

With the sign of \eqref{eq:hhd}, the flux points toward the maxima of the divergence potential $\varphi_\sigma$, which
mark sinks where the agents converge, and away from its minima, which mark sources. The extrema of the vorticity
potential $\psi_\sigma$ are centers of rotation: minima for a counterclockwise rotation and maxima for a clockwise one.
On a bounded domain, the Helmholtz--Hodge decomposition is not unique. It is fixed by boundary conditions, or by setting
aside a harmonic part that stands for the influences from outside the domain \citep{bhatia2014natural}. Here the flux
is a sum of Gaussians, smooth on the whole plane and vanishing far from the agents. Such a field has no harmonic part
\citep[Section~3.1]{bhatia2013helmholtz}, and its two potentials are the solutions that vanish at infinity of the
Poisson equations $\Delta\varphi_\sigma=\nabla\cdot\vVb_\sigma$ and $\Delta\psi_\sigma=\nabla^{\perp}\cdot\vVb_\sigma$
\citep[Eq.~16]{bhatia2013helmholtz}, where $\nabla^{\perp}\cdot\vW=\partial_xW_y-\partial_yW_x$ is the two-dimensional
curl of a vector field $\vW$. The decomposition \eqref{eq:hhd} is then the two-dimensional form of the decomposition
without its harmonic part \citep[Eq.~15]{bhatia2013helmholtz}. A Gaussian is, up to a factor, the Laplacian of a radial
function whose derivative has a closed form, which yields the kernel $\kappa_\sigma$ of \eqref{eq:kappa}
(\cref{app:proofs}); its numerator cancels the singularity of $1/\|\vd\|^2$, so that it equals $1/2$ at $\vd=\vec 0$.
The two potentials \eqref{eq:potentials} rebuild the averaged flux exactly, with no harmonic remainder
(\cref{prop:hhd}). The four fields have no unit, and a replay of the recording at another constant rate, which
multiplies every velocity by the same factor, leaves them unchanged.

The potentials of a Helmholtz--Hodge decomposition also serve to locate the singularities of a vector field
\citep{polthier2003identifying}. They are usually computed on meshes \citep{polthier2003identifying,tong2003discrete},
with particles \citep{petronetto2010meshless} or with matrix-valued kernels fitted to scattered samples
\citep{macedo2010learning}, as reviewed by \citet{bhatia2013helmholtz}; the Gaussian footprint gives them in closed form
instead. \Cref{fig:flows} shows them on three elementary motions.

\begin{figure}[!htbp]
\centering
\figin{\linewidth}{figures/fig_flows.pdf}
\caption{The divergence and vorticity potentials of three elementary motions of $20$ synthetic agents ($\sigma=7$\,cm;
arrows: the averaged flux): a sink is a maximum of $\varphi_\sigma$, a source a minimum, and a counterclockwise
rotation center a minimum of $\psi_\sigma$. The potential that a motion does not produce stays below
$\nm{flows-flat-pct}\,\%$ of the color scale of its row.}
\label{fig:flows}
\end{figure}
\FloatBarrier

\subsection{The coupling}

Taken separately, the diagrams of $\varphi_\sigma$ and $\psi_\sigma$ record how strong a sink and a vortex are, but not
whether the two lie at the same place. Multiparameter persistence would summarize several fields together, but its
summaries do not reduce to a diagram, and comparing them is far less simple
\citep{carlsson2009theory,botnan2023introduction}. The coupling $J_\sigma$ of \eqref{eq:jacobi} is not a new physical
quantity, but it records in one field what the separate diagrams miss. Like the other fields, it is invariant under
rotations and translations, since a rotation does not change a determinant. \Cref{fig:potentials} shows two swarms with
the same ingredients, arranged differently: the diagrams of $\varphi_\sigma$ and $\psi_\sigma$ change only in part, while
that of $J_\sigma$ is empty when the sink and the rotation center coincide and appears when they are apart.


\subsection{Extended persistence and the sliced distance}

Persistent homology follows how the sublevel sets $\{f\leq c\}$ of a field change as the level $c$ rises
(\cref{fig:persistence}). A component is born at each minimum and dies when it merges with an older one. Cycles are born
at saddles and filled at maxima. Ordinary persistence \citep{edelsbrunner2002topological,zomorodian2005computing}
records the level at which each homology class, a connected component or a cycle, is born and the level at which it
dies. Some classes never die, such as the component that holds the global minimum; they are called essential, and their
death would be infinite. Extended persistence \citep{cohen2009extending} completes the ascending sequence of sublevel
sets with the homology of the domain relative to its superlevel sets, swept downward, so that every class has a finite
birth and death. It yields six diagrams per field: the ordinary diagrams $\mathrm{Ord}_0$ and $\mathrm{Ord}_1$, the
relative diagrams $\mathrm{Rel}_1$ and $\mathrm{Rel}_2$, and the extended diagrams $\mathrm{Ext}_0$ and
$\mathrm{Ext}_1$. On a fixed complex, the bottleneck distance between the extended diagrams of two fields is at most the
largest difference of their values \citep[Section~6]{cohen2009extending}. Persistence is multiscale in another sense
than the ladder: at one width, it follows the field across all its levels.

The sliced Wasserstein distance \eqref{eq:sw} is of order one. A field holds a few long-lived structures and many
short-lived ones, whose number and levels also describe the motion, and squared differences would let a single
long-lived structure outweigh all of them. The distance is bounded above by $2\sqrt2$ times the first diagram distance
\citep[Theorem~3.3]{carriere2017sliced}, the least total displacement that matches the points of two diagrams with one
another or with the diagonal, so that a small displacement of the points changes it little. It is differentiable
almost everywhere, since the sort and the absolute value are linear by pieces.

\begin{figure}[!htbp]
\centering
\figin{\linewidth}{figures/fig_persistence.pdf}
\caption{From a field to a distance: (a)~the sublevel sets of a profile; (b)~the resulting points, ordinary and
extended; (c)~two diagrams, each augmented with the projections of the other onto the diagonal (hollow); (d)~their
sorted projections on one direction.}
\label{fig:persistence}
\end{figure}
\FloatBarrier

\subsection{The domain without a grid}

Each recording gets its own seeds at every width, of three kinds. The first kind is a subset of the agent positions,
chosen by farthest-point sampling until every position lies within $\nm{cover-sigma}$ of the subset. The second, for
each field, are the barycenters of the $\nm{frac-pct}\%$ of the Delaunay triangles of these positions where its
gradient is smallest; these seeds lie near the critical points of the field. The third are up to $\nm{n-jacobi}$ points
of the coupling skeleton, located by $\nm{n-bis}$ bisection steps on the Delaunay edges where $J_\sigma$ changes sign.
The five fields are evaluated at the same seeds, and seeds farther than $5\sigma_m$ from every agent at every instant
are discarded. The domain is the alpha complex of the remaining seeds, computed with GUDHI \citep{maria2014gudhi}. It
keeps the simplices of their Delaunay triangulation whose alpha value, that is, whose squared radius, is at most
$(\nm{c-alpha}\,\sigma_m)^2$ \citep{edelsbrunner2010computational}. On this complex, each field $f$ induces a lower-star
filtration: a vertex enters at the value of $f$ there, and an edge or a triangle at the largest value of $f$ at its
vertices.

\subsection{Ranks, fusion, affinity propagation and the modal region}

With mutual ranks, a pair of recordings is close only when each is among the nearest of the other, and the percentile
ranks then give every matrix the same distribution. Weights read from the data, through the spectral gap of each field,
are compared with the weights of \eqref{eq:view} in \cref{app:fusion}. In the fusion \eqref{eq:snf}, the local kernel
keeps distinct behaviors apart, while the diffusion joins the recordings of one behavior that single widths leave split.

Affinity propagation passes two kinds of messages, responsibilities and availabilities, between the recordings. The
messages are updated until a set of exemplars emerges and each recording is assigned to one of them. Its preference,
which sets how readily a recording becomes an exemplar, is the median of the entries of $\Sstar$, its diagonal first set
to its largest value. Its damping and its convergence rule are those of its original implementation
(\cref{tab:settings}), so that the clustering adds no tuned parameter. In the modal consensus, ties for the
exemplar of a group go to the member most often clustered with the others.

\subsection{Splitting and differentiating a distance}

Each term of the sliced distance is shared equally between the two elements it matches, a projection onto the diagonal
standing for the point it projects. Each diagram point then passes its share, half and half, to the two seeds that
realize its birth and its death. Holding the seeds fixed while the trajectories move is exact at first order for the
critical values, whose gradient vanishes. It is not exact for a quantile of the values, which lies at an ordinary point
where the gradient does not vanish (\cref{app:attribution}).

% ============================================================================
\section{Settings and notation}
\label{app:settings}

This appendix lists the settings of the method and tests them one by one, then gathers the notation of the paper.
\Cref{tab:settings} lists the settings, none of which is tuned on the mission labels.

\begin{table}[!htbp]
\centering
\caption{Fixed settings of the method. A \emph{resolution} setting can be refined without changing
what is computed, only how finely; a \emph{design} setting defines what is computed.}
\label{tab:settings}
\tabin{settings}
\end{table}

To check that no setting decides the result, each one was varied alone on Top-1, through the whole method. Each
version was compared without labels with a reference: its finest value for a resolution setting, and the value of
the method for a design choice. The
comparison uses the rank correlation of the distances, field by field and scale by scale, and the agreement
of the taxonomies (\cref{tab:study}). None of them changed the number of behaviors, and at most
\nm{study-runs-moved} of the $250$ runs changed cluster (an adjusted Rand index of at least
\nm{study-ari-min} with the reference). The distances depend mostly on how finely the domain is sampled at
the smallest scale, which is why the seeds follow the scale (\cref{sec:seeds}).

\begin{table}[!htbp]
\centering
\caption{Each setting varied alone on Top-1, without labels. \emph{Distances}: the Spearman rank correlation
between the distances at the value of the method and at the finest value tested, the smallest over
the twenty matrices of fields and scales. For a design choice, it is the smallest over the values tested.
\emph{Taxonomy}: the smallest adjusted Rand index between the consensus partition of a value tested and
that of the reference. Resolution settings are checked on the domain of the method (seeds within
$\sigma/2$ of every robot position) when it was run, otherwise on a fixed subset of $\nm{n-coreset}$ robot positions.}
\label{tab:study}
\tabin{study}
\end{table}

The notation of the paper is gathered in \cref{tab:notation}, grouped by the step of the method that introduces it.

\begin{table}[!htbp]
\centering
\caption{Notation, grouped by the step of the method that introduces it.}
\label{tab:notation}
\begin{tabular*}{\linewidth}{@{\extracolsep{\fill}}l>{\raggedright\arraybackslash}p{6.4cm}l@{}}
\toprule
Symbol & Meaning & Defined in\\
\midrule
\multicolumn{3}{@{}l}{\emph{Recordings and fields}}\\
$N$, $T$, $\Delta t$ & agents, instants per recording, sampling step & \cref{sec:fields}\\
$\vx_i(t)$, $\vv_i(t)$ & position and forward-difference velocity of agent $i$ (vectors) & \cref{sec:fields}\\
$O_i$, $\sigma$ & Gaussian footprint and its width, the observation scale & \eqref{eq:footprint}\\
$M_0$, $\vM_1$, $\tM_2$ & moments of order zero, one and two of the velocity distribution (scalar, vector, tensor) & \eqref{eq:moments}\\
$\vV_\sigma$ & flux, the vector field $\vM_1$ & \eqref{eq:hhd}\\
$\rho_\sigma$, $K_\sigma$ & density and kinetic intensity & \cref{sec:fields}\\
$\varphi_\sigma$, $\psi_\sigma$ & divergence and vorticity potentials of the Helmholtz--Hodge decomposition & \eqref{eq:potentials}\\
$\kappa_\sigma$ & radial kernel of the potentials & \eqref{eq:kappa}\\
$\vbar$, $\overline{v^2}$ & mean speed and mean squared speed of the recording & \cref{sec:fields}\\
\addlinespace
\multicolumn{3}{@{}l}{\emph{Persistence and distances}}\\
$\sigma_1,\dots,\sigma_M$ & ladder of scales, $M=4$ & \eqref{eq:ladder}\\
$J_\sigma$ & Jacobi coupling $\sigma^2\det(\nabla\varphi_\sigma,\nabla\psi_\sigma)$ & \eqref{eq:jacobi}\\
$D^{s}_{f,m}(A)$ & diagram $s$, one of six, of field $f$ at scale $\sigma_m$ for recording $A$ & \eqref{eq:sw}\\
$d_{f,m}$ & sliced Wasserstein distance between recordings & \eqref{eq:sw}\\
\addlinespace
\multicolumn{3}{@{}l}{\emph{Fusion and clustering}}\\
$\tilde d_{f,m}$ & percentile rank of the mutual rank of $d_{f,m}$ & \cref{sec:fusion}\\
$W_m$ & balanced view at scale $\sigma_m$ & \eqref{eq:view}\\
$\knn$, $t$ & neighborhood size and diffusion iterations of the fusion & \eqref{eq:snf}\\
$\Sstar$ & consensus network after fusion & \cref{sec:fusion}\\
$C$, $\Cstar$ & number of clusters, its modal value across the local grid & \cref{sec:ap}\\
\bottomrule
\end{tabular*}
\end{table}

% ============================================================================
\section{The ladder of scales}
\label{app:gallery}

\Cref{fig:scales_all} shows the five fields of the Shelter exemplar of Top-1 at the four scales of
the ladder. At the finest scale the fields resolve individual paths and the walls of the shelter; at
the coarsest they keep the organization of the whole group.

\begin{figure}[!htbp]
\centering
\figin{0.70\linewidth}{figures/fig_scales.pdf}
\caption{The five fields of one recording (the Shelter exemplar of Top-1) at the four scales of the
ladder, $\nm{sigma1-cm}$, $\nm{sigma2-cm}$, $\nm{sigma3-cm}$ and $\nm{sigma4-cm}$\,cm.
Color scales are per panel: sequential for the density $\rho_\sigma$ and the kinetic intensity
$K_\sigma$ (zero is pale blue), diverging for the potentials $\varphi_\sigma$, $\psi_\sigma$ and the
coupling $J_\sigma$ (blue negative, orange positive, zero gray). White marks points farther than
$5\sigma$ from every robot at every instant, where the fields are not evaluated.}
\label{fig:scales_all}
\end{figure}
\FloatBarrier

The top of the ladder follows from the merging width $D_s/2.35$ at which the footprints of the whole
group stop being resolved (\cref{sec:ladder}). \Cref{fig:uncapped} shows what lies near it. On a ladder
extended up to that width, the density and the vorticity potential of the Homing exemplar keep their
organization up to a few tens of centimeters. Past that, each field becomes one smooth bump over the
whole group, and its diagrams no longer depend on how the robots moved. The cap of the method therefore lies far
below the merging width. $D_s$ measures the extent of the swarm, not the size of the structures of its
motion, such as the corridor, the shelter or a group of robots.

\begin{figure}[!htbp]
\centering
\figin{\linewidth}{figures/figA_uncapped.pdf}
\caption{The fields near the merging width. The density and the vorticity potential of the Homing
exemplar on a ladder that extends to the merging width of the whole group, $D_s/2.35$. Beyond a few
tens of centimeters every field becomes a single featureless bump, whose diagrams no longer
distinguish behaviors.}
\label{fig:uncapped}
\end{figure}
\FloatBarrier

The method uses the cap $k=\nm{cap-k}$, and its taxonomy hardly depends on this value. For caps
$\sigma_{\max}=D_s/(2.35\,k)$ from a third of the merging width down to a twentieth of it, the whole method was run on
Top-1 and the taxonomy obtained as in \cref{sec:ap}. Every cap from $k=\nm{cap-same-kmin}$ to $k=\nm{cap-same-kmax}$, a
top of the ladder between $\nm{cap-same-cmmin}$ and $\nm{cap-same-cmmax}$\,cm, gives the same taxonomy, and
$k=\nm{cap-apart-k}$ moves \nm{cap-apart-runs} run. Even with no cap at all, on a fixed subset of $\nm{n-coreset}$ robot
positions, the taxonomy differs from that of the capped ladder by \nm{uncapped-runs-moved} run, and adding scales between
the same bounds, up to $\nm{uncapped-mmax}$, leaves it unchanged.

% ============================================================================
\section{Robustness of the fusion}
\label{app:fusion}

\subsection{Alternative fusion architectures}
\Cref{tab:designs} compares, on Top-1 and over the same $99$ settings, our fusion with
alternatives built from the same rank-normalized matrices and clustered by the same affinity
propagation. Without diffusion, affinity propagation splits the runs of each single controller of
Top-1 into several groups, $\nm{design-none-C}$ in all, although each mission is one controller there,
and every setting returns the same partition. Diffusing first among the
fields of each moment, then across scales, then across moments, compounds the smoothing. The pair $(\knn,t)$
no longer means a single resolution, missions merge already at moderate settings (median purity
\nm{design-hierarchical-medpur}), and the number of clusters stops decreasing regularly with $t$.
A single SNF
over all field-and-scale views gives the flux three times the weight of the density. Among the
architectures that diffuse, the two fusions that balance the moments and the single SNF over twelve views, one per
moment and scale, reach a median purity of at least \nm{design-top-medpur-min}, and the two others of at most
\nm{design-low-medpur-max}. The settings in
which purity and ARI are both one number \nm{design-balanced-perfect} for our fusion, \nm{design-fluxrerank-perfect}
for the fusion that re-ranks the flux first, and \nm{design-snf12-perfect} for the single SNF over twelve views.
Our fusion keeps the number of clusters from increasing with $t$ at \nm{design-balanced-mono} of the nine
neighborhood sizes, the re-ranked fusion at \nm{design-fluxrerank-mono}, and the single SNF over twelve views at
\nm{design-snf12-mono}.

\begin{table}[!htbp]
\centering
\caption{Fusion architectures on Top-1 (same distances, same clustering, $99$ settings).
\emph{Modal}: settings at the modal count $\Cstar$; \emph{perfect}: settings with purity and ARI equal to
one; $C\le2$: settings where the partition collapses; \emph{monotone}: neighborhood sizes for which
$C$ never increases with $t$.}
\label{tab:designs}
\tabin{designs}
\end{table}


\subsection{Rank fusion against standard SNF}
On Top-1, with $\knn=5$, rank fusion settles on \nm{std-ours-C} clusters from $t=\nm{std-ours-t-from}$ on and
keeps them up to $t=\nm{std-t-max}$. With the affinity of standard SNF on the raw distances instead of ranks, the
number of clusters drifts with the number of diffusion iterations, from \nm{std-C-first} at
$t=\nm{std-t-first}$ to \nm{std-merge-C} from $t=\nm{std-merge-from}$ on, where two missions share a cluster
(purity \nm{std-merge-purity}).

\subsection{Convergence of affinity propagation}
With damping $0.9$ the messages of affinity propagation change slowly, so the algorithm must be
given enough iterations to settle. The settings of the original implementation declare convergence
after a window of $100$ unchanged iterations, within at most $1000$. With them it converges at every
setting of Top-1 and of Top-10, the largest dataset, after a median of \nm{ap-top10-iter-median}
iterations on Top-10 and at most \nm{ap-top10-iter-max}. A window of $200$ or $400$ iterations, or a ceiling of $5000$, leaves
each of the $99$ partitions of both datasets unchanged: the algorithm has time to finish. The
shorter window of scikit-learn, $15$ iterations, stops earlier and changes \nm{ap-short-top1-changed}
partitions of Top-1 and \nm{ap-short-top10-changed} of Top-10. Another damping, $0.8$ or $0.95$, or
another draw of the tie-breaking noise can reach another set of exemplars, at up to
\nm{ap-other-changed-max} of the $99$ settings. The modal consensus taxonomy stays the same,
with an adjusted Rand index of at least \nm{ap-other-cons-min} with the reference.

\subsection{Each step of the fusion changed alone}
\begin{table}[!htbp]
\centering
\caption{Each step of the fusion changed alone, on the ten datasets, the taxonomy of every variant
built on its own modal region, as for the method. \emph{Modal} and \emph{stability}: means over the datasets, without
labels. \emph{Purity} and \emph{ARI}: means over the datasets. \emph{Misassigned}: runs outside a
cluster dominated by their mission, summed over the datasets.}
\label{tab:variants}
\tabin{variants}
\end{table}

\Cref{tab:variants} changes one step of the fusion at a time and keeps the others. On the ten
datasets, the taxonomy of every variant is built on its own modal region, as for the method. The first step is the affinity
given to each view. Percentile ranks alone, without the mutual ranks, misassign
\nm{fx-affpercentile-wrong} runs, \nm{fx-affpercentile-xs} of them XOR Aggregation runs placed with
Shelter (\cref{app:merge}). Percentile ranks passed through the normal quantile function misassign
\nm{fx-affprobit-wrong}. A Gaussian kernel whose width adapts to each pair, through the distances of the two runs to their seventh neighbors, keeps the purity. It splits the missions
further, however, with a mean adjusted Rand index of \nm{fx-affselftuning-ari}. Inside one scale, an equal weight per field misassigns
\nm{fx-fieldsequal-wrong} runs. The weights can also be read from the data. The spectral gap of a
field at one scale is the largest gap between two consecutive eigenvalues of the normalized Laplacian of its
neighbor graph, the first gap left out \citep{von2007tutorial}. This graph joins two runs when one is among the
seven nearest neighbors of the other, with the weight $e^{-d_{f,m}/s}$, where $s$ is the standard deviation of
the entries of $d_{f,m}$. A larger gap means that the field splits the runs into clearer groups. Given to each
field alone, these gaps misassign \nm{fx-fieldsconfidence-wrong} runs. Shared within each moment, that is,
divided by the number of fields of the moment and normalized to sum to one, they give the weights read from the
data that \cref{sec:fusion} mentions. They reduce to the weights of \eqref{eq:view} when the gaps are equal, and they misassign
\nm{fx-fieldssourceconf-wrong} runs. The
weighted geometric mean misassigns \nm{fx-fieldsgeomean-wrong} runs, and the maximum over the fields
\nm{fx-fieldsmax-wrong}. Across scales, the diffusion is needed: averaging the four views without it returns one partition at every setting, which splits the
missions into \nm{fx-mean-cmin} to \nm{fx-mean-cmax} groups as controllers are pooled. A single SNF over
the twenty field-and-scale views merges missions, with \nm{fx-scalessnf20-wrong} runs misassigned. On
the consensus, affinity propagation with the minimum preference reproduces the taxonomy of the method, with the
same $\Cstar$ on every dataset and an adjusted Rand index of at least \nm{fx-apmin-ari-min} with it. Louvain
modularity finds the same $\Cstar$ on \nm{fx-louvain-same} of the ten datasets, at an adjusted Rand index of at least
\nm{fx-louvain-ari-min}. Clusterings that choose the number of groups from an
eigengap, from the largest gap of an average-linkage tree or from density merge missions instead,
into \nm{fx-collapse-cmin} to \nm{fx-collapse-cmax} groups, with a purity of at most
\nm{fx-collapse-pmax}. A higher adjusted Rand
index in the table goes with fewer groups on the pooled datasets, not with fewer misassigned runs,
since that index counts the sub-behaviors against the method (\cref{sec:protocol}).


% ============================================================================
\section{Other descriptions on all the runs}
\label{app:baselines}

\Cref{tab:baselines_full} compares the descriptions of \cref{tab:baselines} on Top-1 and on Top-10, where the missions
split into sub-behaviors. There, homogeneity and completeness are the scores to read, since the ARI counts the
sub-behaviors as errors.

\begin{table}[!htbp]
\centering
\caption{Descriptions of the runs under the same fusion, affinity propagation and modal consensus, on Top-1 and Top-10. The
number of modal settings follows $\Cstar$ in parentheses; \emph{error-free}: settings, out of $99$, in which no run is
misassigned, those in the modal region in parentheses. The method is shaded. $^{*}$Partition published for the same Top-1
runs, with five fixed clusters \citep{legarda2026universal}.}
\label{tab:baselines_full}
{\footnotesize\renewcommand{\arraystretch}{1.0}\tabin{baselines}}
\end{table}
\FloatBarrier

% ============================================================================
\section{Rotations on the robot data}
\label{app:invariance}

To test the invariance on the robot data, five runs of Top-1, one per mission, were rotated by random angles about the center of the arena
and processed again with the same code. For every field and scale, the distance between each run and
its rotated copy was divided by the median distance between the run and twelve other runs of its mission. The
density, the kinetic intensity and both potentials are recovered up to rounding, within \nm{inv-other-max} of that
median distance. The coupling differs most at the finest scale, by at most \nm{inv-worst} of it, because it is the
most sensitive to the placement of the seeds.

% ============================================================================
\section{XOR Aggregation runs that resemble Shelter}
\label{app:merge}

\begin{figure}[!htbp]
\centering
\figin{\linewidth}{figures/figA_merge.pdf}
\caption{XOR Aggregation runs that percentile ranks alone, without mutual ranks, place in the Shelter cluster on
Top-4. (a)~Median fraction of robots
moving faster than $1$\,cm/s over the run, for these runs, for the other XOR runs and for the Shelter
runs. (b--d)~Final positions of the robots, pooled over the runs of each group.}
\label{fig:merge}
\end{figure}

With percentile ranks alone instead of mutual ranks, XOR Aggregation runs join the Shelter cluster. On Top-3, all
\nm{merge-top3-n} XOR Aggregation runs do, and on Top-4, \nm{merge-top4-n} of the $200$ do, from every controller of the
mission. The method places none of them with Shelter on either dataset. Compared with the other XOR runs of Top-4, these
runs settle late, after a median of \nm{merge-top4-settle-abs}\,s against \nm{merge-top4-settle-kept}\,s. They also move
about twice as fast on average (\nm{merge-top4-speed-abs} against \nm{merge-top4-speed-kept}\,m/s, and
\nm{merge-top4-speed-shel}\,m/s for Shelter), yet they end on the same two black areas (\cref{fig:merge}). What keeps them
with their mission is their neighborhood. Among their seven nearest neighbors in each view, Shelter runs make up
\nm{merge-top4-nn-pct}\,\% with percentile ranks, against \nm{merge-top4-nn-mut}\,\% with mutual ranks. The Shelter runs
have one another as nearest neighbors, so their mutual ranks with these runs are large.

% ============================================================================
\section{Stability and reproducibility of the attribution}
\label{app:attribution}

The attribution holds the seeds fixed, which is justified at the critical points inside the domain.
There the gradient of the field vanishes, so that a critical value does not change at first order when
its point moves with the trajectories, as the envelope theorem states. On a synthetic milling group,
perturbing the tracks changes the minimum of $\psi_\sigma$ by the same amount, to the precision of the
finite differences, whether the minimum is relocated or held in place. At a seed that
moves with a robot, as the seeds taken from the robot positions do, the two derivatives differ even in sign
(\nolinkurl{manuscript/verify_propositions.py}). The value distance splits in the same way as the sliced
distance: the term of each quantile is shared equally between the two runs, and the half of a run between the two
order statistics that the quantile interpolates.

The split of \cref{sec:res_attribution} was checked in two ways. First, a Gaussian noise was added to every
position of the Directional Gate exemplar, with the seeds held fixed. A noise of $2$\,mm changes the field at the
seeds by up to \nm{attr-stab-field-change} of its largest value. It leaves \nm{attr-stab-pers} of the topological
split and \nm{attr-stab-val} of the value split on the same seeds, and the topological split is the more stable at
every noise tested (\nm{attr-stab5-pers} against \nm{attr-stab5-val} at $5$\,mm). The term of a quantile falls on
the seed that holds that rank, which a small noise can change, while the share of a diagram point stays on its
critical points. Second, the split was repeated over the \nm{attr-runs} runs of the Directional Gate controller,
each against the same Foraging exemplar. For the persistence diagrams, the two extrema with the largest shares are
a maximum and a minimum of $\psi_\sigma$, the two rotation centers of the traffic, \nm{attr-centres-phrase}
\nm{attr-runs} runs, with a median share of \nm{attr-centres-share}. Half of the mean map lies in
\nm{attr-mean-half-pers} of the arena, against \nm{attr-mean-half-val} for the value distributions, whose split
covers the whole positive cell of the potential. The centers themselves move by about \nm{attr-centres-spread} from
run to run, so two single topological maps overlap less on average (\nm{attr-overlap-pers}) than two diffuse value
maps (\nm{attr-overlap-val}). What recurs from run to run is the structure, not its exact position.

The split for the ten pairs of missions is given in \cref{tab:attribution} (\cref{sec:res_attribution}). In
\nm{attrp-pairs-50} of the ten pairs, the two largest extrema are a maximum and a minimum in every one of the
\nm{attr-runs} runs; only Directional Gate against XOR Aggregation drops to $\nm{attrp-recur-min}/50$, because one run
splits its largest weights across two maxima of $\rho_\sigma$ before reaching its minimum.



\end{appendices}
\fi

\clearpage
\bibliography{references}

\end{document}
